% ==================== 第1章 ====================
\chapter{绪论}
\label{ch:intro}

\section{研究背景与问题提出}

金融分析中的问题并不都要求复述已经披露的事实。证券分析师阅读年报时，既会询问已经公告的营业收入，也会询问下一年度营业收入能否延续增长、毛利率能否维持、业绩指引能否兑现，以及年报将在何时披露。前一类问题的事实已经存在于截止时点之前的申报文件、公告或结构化数据中，回答的重点是定位证据并完成计算。后一类问题的真值在截止时点尚未实现，系统只能依据当时可得的证据形成一个可修正的判断。两类问题共享同一套金融文本，却要求不同的证据表示和输出接口。

现有金融问答研究主要服务第一类问题。FinQA研究基于财务报表的数值推理\citep{chen2021finqa}，TAT-QA研究表格与文本的混合推理\citep{zhu2021tatqa}，FinanceBench把问答扩展到真实金融文档\citep{islam2023financebench}，DocFinQA进一步把证据范围从段落扩展到完整10-K\citep{reddy2024docfinqa}。这些基准的答案在文档中已经存在，评价指标是精确匹配或词级F1。它们衡量阅读器能否找到并计算答案，不衡量系统能否判断一条断言是否应当以点答案回应。

金融断言验证研究把给定陈述映射为支持、反驳或证据不足。FinDVer面向长文档和混合内容的可解释断言验证\citep{zhao2024findver}，Fin-Fact面向多模态金融事实核查\citep{rangapur2025finfact}。这类任务承认证据可能不足，却通常假设断言已经给定、证据已经检索，输出是一个离散裁决。对于“下一年度营业收入增速会落在哪个区间”这样的问题，验证器的标签空间并不适用，因为结果尚未实现，证据只能支持某个区间，而不能证明该区间已经成为事实。\par
时点检索和预测基准进一步说明了这一区别。TempFinRAG在检索前执行as-of日期边界，把时间证据泄漏压到较低水平\citep{tao2026tempfinrag}；EPSBench在财报发布前预测EPS，发现多数LLM并不显著优于分析师一致预期，名义80\%区间的实际覆盖率只有37\%至70\%\citep{yin2026epsbench}；FinVerify对金融agent的数值断言做时间完整性、来源保真和会计一致性检查，并允许弃权\citep{guan2026finverify}。FinAbstain进一步把时点检索、多模态证据、证据核验、不确定性校准和显式弃权组合成选择性金融预测流程\citep{torres2026finabstain}。这些工作覆盖了时点有效性、数值预测、证据核验和弃权动作，但FinAbstain的证据核验仍是不确定性聚合之前的过滤，尚未把多类别资格判定作为独立阶段并与类型可靠性演化耦合。

检索增强生成（RAG）把检索器与生成器组合，假设答案可以由检索到的上下文拼装\citep{lewis2020rag}。Self-RAG通过自我反思决定是否检索\citep{asai2024selfrag}，CRAG在检索质量不足时触发纠正性检索\citep{yan2024crag}，FLARE在生成过程中按置信度触发检索\citep{jiang2023flare}，Adaptive-RAG按问题复杂度选择检索策略\citep{jeong2024adaptiverag}，智能体框架交错进行推理、检索和自我修正\citep{yao2023react,shinn2023reflexion}。这些方法改善了证据获取，却大多以一个生成答案作为输出。当前瞻性断言的结果尚未实现时，单一答案无法表达证据不足、结果不确定和错误成本不对称。

管理层预测研究提供了另一个视角。业绩指引同时包含区间、时点和可验证性，不能等同于已实现业绩\citep{hirst2008management}。指引的宽度、管理层的历史兑现记录、当前季度相对指引区间的位置、行业需求和产能约束，都会影响最终结果。把这些因素压缩为一个类别标签，会丢失证据的作用域、来源和时效；把它们留给一个端到端模型，又会失去可审计性。金融分析需要的是一种中间状态，使证据到判断、判断到决策的每一步都可以被检查。

本文关注前瞻性金融断言的裁决问题。给定一条断言、一个截止时点和一个有限结果空间，系统只能使用截止时点之前的证据，输出一个点答案、一个裁决集或弃权。本文把这个任务形式化为断言契约，提出FinPlanRAG框架，其核心是AnchorEvolve锚点演化模块。契约明确公司、目标期间、指标、结果空间和截止时点；AnchorEvolve从时点证据中抽取带来源的锚点，对锚点做空泛、过期、对抗、硬、软五类资格判定；已揭示结果按公开揭示日更新锚点类型的可靠性；校准后验在点答案、非单元素裁决集和弃权三个动作上最小化期望损失。本文不试图用一个模型替代分析师，而是把分析师的证据判断过程写成可检查、可更新、可复核的状态。

\section{研究问题}

本文的研究问题可以分解为四个相互衔接的问题。

第一，如何表示前瞻性断言的证据。已有系统通常把证据表示为检索到的文本片段或一个分类标签。文本片段缺少结构，分类标签缺少来源。前瞻性断言需要一种表示，既保留证据的语义关系，又记录证据的作用域、来源、时间和可靠性。本文用三层锚点表示回答这个问题：稳定关系层定义可复用的关系类型，行业模式层把关系绑定到具体字段和生效规则，实例层记录每条证据的取值、来源和资格。

第二，如何避免把未来信息泄漏到时点判断。前瞻性断言的真值在截止时点之后才揭示。如果可靠性更新、先验估计或权重选择使用了未来结果，系统在回测中会表现良好，在部署中却不可复现。本文用公开揭示日约束状态更新，只有揭示日早于当前截止时点的结果才能进入可靠性后验和先验；训练、校准和权重选择只在开发集与验证集上进行。

第三，如何把语义判断转换为风险可控的决策。LLM可以从文本中抽取锚点、解释证据和形成语义判断，但其概率输出往往存在类别偏置和校准误差。统计模型可以使用结构化特征进行概率校准，却缺少对文本证据的理解。本文让LLM构造锚点状态，让校准随机森林提供结构化后验，再在同一个损失函数下选择点答案、裁决集或弃权。语义与校准的分工由验证集选择，而不是由测试集比较后择优。

第四，如何评价这套框架。公开基准中，DocFinQA、FinQA、TAT-QA和FinanceBench的答案已经存在于文档中，FinDVer和Fin-Fact的断言已经给定、标签是支持、反驳或证据不足。没有一个公开基准同时包含断言、截止时点、结果空间、时点证据和已实现结果五元组。本文因此构建收入指引断言基准，并明确说明公开基准的适配边界：它们可以检验证据检索、阅读器和断言资格判定等组件，但不能直接检验前瞻性结果预测。本文不把组件级结果包装成端到端优势。

\section{研究内容与方法}

本文的研究内容包括四个部分。

第一，断言契约与问题形式化。本文将前瞻性断言表示为一个五元组，明确结果空间、截止时点和可用证据集合。结果空间可以是数值指标的三个区间，也可以是披露日期的三个日历窗口。契约把“回答什么”和“何时回答”写成显式对象，使不同方法可以在同一口径下比较。

第二，FinPlanRAG框架。FinPlanRAG以AnchorEvolve为核心模块，框架由证据构造、锚点抽取、资格判定、时间演化、信念更新和全动作读出六个模块组成。证据构造从公司公告、监管文件和网络检索中获取带日期的材料，按截止时点过滤并分配来源编号。锚点抽取调用LLM，要求每条锚点引用来源编号、给出逐字引文、发布日期、隐含结果标签和校准可靠性。资格判定用确定性规则检查空泛、过期、对抗、硬、软五类状态。时间演化按锚点类型维护Beta可靠性，只用已公开揭示的结果更新。信念更新把软锚点转换为似然，把硬锚点转换为约束，再与LLM的标签分数和随机森林后验融合。全动作读出在点答案、非单元素裁决集和弃权上计算期望损失，选择损失最小的动作。

第三，实证评估。本文在1,005条季度营业收入指引断言上评估框架，其中227条测试断言具有XBRL核验或XBRL覆盖标签。数据按公司不重叠、按年份训练的方式划分，使用五个随机森林种子。本文比较纯LLM锚点后验、锚点与随机森林的融合后验、随机森林诊断和六个检索与验证基线，报告准确率、平衡准确率、宏F1、决策损失、可信集覆盖率和平均集合大小。本文还报告文本锚点、时间演化和全动作读出的消融，并给出置信区间和McNemar检验。

第四，可审计接口。本文实现一个本地命令行工具和Claude Code的MCP接口。用户可以用自然语言提出前瞻性断言，系统完成证据检索、锚点构造、披露日期预测和损失感知裁决，并把每条证据、每个锚点、每次信念变化和每个动作的期望损失写入审计文件。本文以中国A股中际旭创2023和2024年年报为例，展示这套接口在历史时点上的运行过程与结果验证。

研究方法上，本文把LLM用于语义抽取，把确定性规则用于证据资格，把Beta后验用于时间演化，把随机森林和温度缩放用于概率校准，把期望损失最小化用于动作选择。每个模块都有明确的输入输出和可检查的中间状态，模块之间通过结构化对象连接，而不是通过一个端到端模型的黑箱输出连接。

\section{主要贡献}

本文的贡献包括六个方面。

第一，本文将前瞻性金融断言形式化为带截止时点和有限结果空间的契约，并把输出从单一答案扩展为点答案、非单元素裁决集和弃权。这个形式化使“证据不足时应当输出什么”成为契约的一部分，而不是事后附加的阈值规则。

第二，本文提出三层锚点状态和五类资格判定。锚点不是检索片段，也不是分类标签，而是带有关系类型、作用域、来源、时间、可靠性和隐含结果的结构化证据。资格判定把空泛、过期、对抗、硬、软分开处理，使每个被拒绝的锚点都有明确理由，每个被采纳的锚点都可以追溯。

第三，本文给出时点正确的时间演化机制。锚点类型的Beta可靠性只使用公开揭示日早于当前截止时点的结果更新，先验只使用已揭示结果构造。这个约束使可靠性演化和先验估计在回测与部署中具有相同的含义。

第四，本文给出全动作损失感知读出，并在三个基座上验证它相对受限可信集门的改进。读出把点答案、裁决集和弃权放在同一个期望损失下比较，使系统可以在不确定性高时选择集合或弃权，而不是被迫输出一个点答案。本文同时说明，可信集是覆盖口径，裁决集是决策口径，两者不应混同。

第五，本文给出一个可运行的命令行与Claude Code接口，把论文中的流程落到真实未决断言上。接口保留了证据来源、原文引文、资格判定、可靠性、信念轨迹和动作损失，为金融分析中的复核和问责提供了可复用格式。\par
第六，本文将类型库扩展纳入算法训练流程。模型在验证误差上提出新的锚点类型和抽取规则，只有通过选择集损失检验的类型才进入类型库；冻结后的类型库作为算法输入，测试集只评估一次。类型库扩展因此不是独立模块，而是与融合权重、温度缩放并列的验证集选择步骤。

\section{结构安排}

本文余下部分安排如下。第\ref{ch:related}章梳理金融问答、断言验证、检索增强生成、选择性预测和时点证据的相关研究，归纳研究缺口。第\ref{ch:problem}章给出断言契约、结果空间、损失函数、研究问题和研究假设。第\ref{ch:framework}章介绍FinPlanRAG的总体框架和AnchorEvolve的三层锚点状态。第\ref{ch:algorithm}章给出完整算法设计，包括证据构造、锚点抽取、资格判定、时间演化、类型库扩展、信念更新、全动作读出和披露日期断言。第\ref{ch:experiment}章报告主基准、A股回测和FinDVer组件检查的实验结果。第\ref{ch:ablation}章报告资格判定、融合权重、类型库扩展和数值桥的消融。第\ref{ch:conclusion}章讨论结果的含义、局限和结论。案例研究放在附录\ref{app:case}。

% ==================== 第2章 ====================
\chapter{文献综述}
\label{ch:related}

本章围绕金融问答与断言验证、检索增强生成与自适应检索、选择性预测与可信集、部分识别与时点证据四条线索梳理相关研究，并在此基础上归纳本文要处理的研究缺口。

\section{金融问答与断言验证}

金融问答基准测试模型在申报文件、表格和报告上的数值推理能力。FinQA要求模型在财务报表中定位数字并完成多步计算\citep{chen2021finqa}，TAT-QA要求模型联合阅读表格和文本\citep{zhu2021tatqa}，FinanceBench把问题扩展到真实公司的公开文件\citep{islam2023financebench}，DocFinQA把证据范围从段落扩展到完整10-K\citep{reddy2024docfinqa}，FinBen提供覆盖多种金融任务的综合评测\citep{xie2024finben}。这些基准的共同点是答案在文档中已经存在，评价指标主要是精确匹配、词级F1或数值容差。它们衡量阅读器能否找到并计算答案，不衡量系统能否判断一条断言是否应当以点答案回应。

金融申报文件的结构与通用问答语料不同。10-K和10-Q通常包含多级标题、跨页表格、脚注和附注，同一指标可能在正文、表格和管理层讨论中以不同口径出现；8-K按事项编号披露业绩、指引、并购和监管事件。XBRL数据提供结构化财务事实，但季度口径、单位、重述和可用日期仍需与文本对齐。通用阅读器擅长语义匹配，却不一定能处理期间、单位和证据闭合。金融问答因此需要把文档结构、时点约束和数值计算纳入同一证据策略。

断言验证基准测试模型判断一条陈述是否被证据支持、反驳或无法判断。FinDVer面向金融领域的断言验证，要求模型在长文档和混合内容中给出裁决和解释\citep{zhao2024findver}；Fin-Fact面向证据驱动的金融事实核查，强调多模态证据和解释生成\citep{rangapur2025finfact}。这类任务的输出是离散裁决，评价指标是裁决准确率和解释质量。验证器通常假设陈述已经被给定，且证据已经被检索，因而没有回答“这条陈述是否应当进入验证流程”的问题。

把验证器直接用于前瞻性断言会遇到标签空间不匹配。验证器的标签是支持、反驳和证据不足，前瞻性断言需要的是结果区间的概率或集合。验证器可以判断“公司说营收将增长”是否有证据支持，却不能判断“公司实际营收将落在指引区间内”的概率。前者是文本蕴含问题，后者是结果预测问题。本文的锚点状态保留了验证器的证据裁决思路，但把标签空间换成结果桶，把证据资格与结果信念分开。

\section{检索增强生成与自适应检索}

RAG将检索器与生成器组合，假设答案可以由检索到的上下文拼装\citep{lewis2020rag}。主动检索和自适应检索进一步决定何时检索或检索多少上下文\citep{jiang2023flare,jeong2024adaptiverag}，智能体框架则交错进行推理、检索和自我修正\citep{yao2023react,shinn2023reflexion}。Self-RAG通过自我反思决定是否检索和如何评价证据\citep{asai2024selfrag}，CRAG在检索质量不足时触发纠正性检索\citep{yan2024crag}，FLARE在生成过程中按置信度触发检索\citep{jiang2023flare}，Adaptive-RAG按问题复杂度选择检索策略\citep{jeong2024adaptiverag}。

这些方法改善了检索过程和证据使用，但输出通常仍是生成的答案。它们可以决定“是否检索”和“检索多少”，却不能决定“返回答案、裁决集还是弃权”。在确定信息问题上，这种输出是合适的；在前瞻性断言上，单一生成答案无法表达结果分布、证据作用域和决策风险。检索策略的自适应并不等于输出契约的自适应，二者是两个不同层次的决策。

本文借用自适应检索的迭代思想，但把停止条件从“检索轮数”改为“证据资格与信念稳定”。系统在证据不足时可以继续检索，也可以在损失允许时输出裁决集或弃权。检索的终点不是固定预算，而是证据状态是否足以支持一个低损失动作。

\section{预测基准、时点泄漏与证据充分性}

与本文最接近的另一条线索是预测基准和时点泄漏。ForecastBench用动态生成的未来问题评估模型预测能力，问题在提交时没有答案，从而降低静态基准的训练污染\citep{karger2025forecastbench}。EPSBench让LLM agent在财报发布前预测季度EPS，发现多数模型并不显著优于分析师一致预期，名义80\%区间的实际覆盖率只有37\%至70\%\citep{yin2026epsbench}。EPSBench的输出是点预测和预测区间，损失是预测误差和区间覆盖；本文的契约输出是点答案、非单元素裁决集或弃权，损失是不对称错误成本下的期望损失。两者在决策问题的形式上不同。这提醒本文，强分析师一致预期和结构化模型是不能省略的基线，区间校准也不能只看名义覆盖率。

TempFinRAG在检索前执行SEC filing可用性边界，把文本、表格、图像和XBRL事实按as-of日期对齐，并报告时间证据泄漏从10.8\%降到1.7\%\citep{tao2026tempfinrag}。ChronoBERT和ChronoGPT用时间受限语料训练，从源头减少前视偏差\citep{he2026chrono}。这些工作说明，时点约束必须进入系统设计，而不是事后过滤；但它们的评价仍以问答准确率或预测误差为主，不处理证据类型可靠性的演化，也不把点答案、裁决集和弃权放进同一个损失。AnE用“anchor”指训练数据的真值锚定，本文的锚点指证据状态中的结构化对象，二者机制不同\citep{wang2026ane}。

证据充分性和选择性预测提供了另一组工具。Sufficient Context把“检索上下文是否足以回答”作为一个独立判断，并用它引导弃权\citep{joren2025sufficient}。FinVerify面向金融agent的数值断言，检查时间完整性、来源保真、可比性和会计一致性，输出接受、拒绝或弃权\citep{guan2026finverify}。SelectLLM把选择性预测放进训练目标，平衡覆盖率和风险\citep{mao2025selectllm}。这些方法能判断“证据够不够”或“该不该弃权”，但通常不维护证据类型的时间可靠性，也不把证据资格判定作为结果信念的一部分。

上述工作共同覆盖了时点检索、预测校准、证据充分性和选择性输出。FinAbstain把其中多项组合成一个选择性金融预测流程，是本文最接近的工作\citep{torres2026finabstain}。本文的缺口因此更窄：不是缺少一个选择性预测框架，而是缺少把证据资格判定作为独立前置阶段、并让资格结果与证据类型可靠性演化耦合的框架和基准。FinAbstain的证据核验在不确定性聚合前过滤明显不可用的证据；本文的五类判定区分空泛、过期、对抗、硬和软五种失效模式，并让被接受的证据进入按揭示日更新的Beta可靠性。这个差别是结构粒度上的，不是决策动作空间上的。

\section{选择性预测、校准与可信集}

拒识选项在识别错误和拒识率之间进行权衡\citep{chow1970optimum}，选择性分类通过置信度阈值在低置信度输入上弃权\citep{geifman2017selective}。现代分类器通常存在校准误差，温度缩放是常用基线\citep{guo2017calibration}；代价敏感学习为不同错误赋予不同价格\citep{elkan2001cost}；分裂共形预测把得分转换为在可交换性下具有分布无关覆盖率的集合\citep{vovk2005algorithmic,angelopoulos2021gentle}。分类共形方法进一步构造自适应集合\citep{romano2020classification}，分布漂移下的自适应共形方法允许覆盖率随时间调整\citep{gibbs2021adaptive}。

这些方法提供了损失、置信度和集合构造的工具，但通常假设任务和回答接口已经固定。选择性分类器在“答案”和“弃权”之间选择，验证器在“支持”“反驳”和“证据不足”之间选择，共形方法在给定标签空间上构造集合。它们没有解决“标签空间本身是否适用于当前查询”的问题。对于前瞻性断言，标签空间是结果桶，证据是锚点，输出还需要说明证据的作用域和时效；这些结构不在标准选择性预测的范围内。

本文借用选择性预测的损失框架和共形集合的覆盖思想，但把标签空间、证据资格和输出动作放在同一个锚点状态上。校准随机森林和温度缩放用于控制后验概率，损失感知读出用于在答案、裁决集和弃权之间选择。与标准共形方法相比，本文的可信集由后验质量构造，覆盖率在实验中被报告，但并未假设交换性。金融时点数据的分布会随时间变化，交换性假设很难成立，因此本文把可信集作为审计口径，把损失最优裁决集作为决策口径。

\section{部分识别、锚点与时点证据}

部分识别表示与现有证据一致的结果集合\citep{manski2003partial}。在金融机器学习中，模型决策只能使用当时可得的信息，否则会产生前视偏差\citep{lopezdeprado2018advances}。本文借用约束和可行集的思想，将锚点定义为来源与假设之间的关系：硬锚点可以排除某个结果桶，软锚点只重新加权结果信念；证据过滤和模式过滤均按截止时点执行。

与部分识别相比，本文的锚点不是完整的结构模型，而是由LLM从文本中抽取的语义约束。锚点带有类型、取值、作用域、来源、时间、可靠性和有效期，因而可以被逐条检查。部分识别给出的是在给定模型下与证据一致的参数集合，锚点给出的是在给定断言下与证据一致的结果桶约束。两者都强调“证据不足时不要强行点识别”，但本文更关注证据的可审计性和可更新性。

时点约束贯穿本文的证据构造、锚点抽取、可靠性更新、先验估计和权重选择。证据只接受可用日期不晚于截止时点的文档，锚点只接受截止时点之前的来源，随机森林只使用截止时点之前可得的特征，校准和权重选择只在开发集与验证集上进行。这一约束使“当时可得信息”成为系统设计的一部分，而不是事后评价的附加条件。

\section{研究缺口}

现有研究分别推进了长文档检索、金融问答、断言验证和低置信度拒识。FinAbstain已经把时点检索、多模态证据、证据核验、不确定性校准和显式弃权组合成一个选择性金融预测流程\citep{torres2026finabstain}。本文的缺口不在这些组件的组合，而在证据资格判定的结构粒度：把资格判定作为先于信念更新的独立多类别阶段，并让被接受的证据按类型和公开揭示日更新可靠性。本文认为，这一缺口可以分解为三个相互衔接的挑战。

第一个挑战是证据表示。现有验证器把断言直接映射为支持、反驳或证据不足，阅读器则直接生成答案。两者都没有保留证据的作用域、可靠性、时效和来源，也没有在证据变化时更新判断。对于前瞻性断言，证据往往只支持某个结果桶、只在某个时间范围内有效、只对某个公司或行业成立；把这些信息压缩为一个标签，系统就无法解释判断为何改变，也无法在结果揭示后更新可靠性。

第二个挑战是时间演化。前瞻性断言的真值在截止时点之后才揭示。如果可靠性更新、先验估计或权重选择使用了未来结果，系统在回测中会表现良好，在部署中却不可复现。现有研究很少把揭示日写进状态更新规则，也很少区分“锚点在截止时点是否可得”和“锚点在结果揭示后是否正确”。这两个问题需要分开处理。

第三个挑战是决策口径。普通准确率把三类错误视为同价，而金融分析对过度乐观、漏报和集合宽度的代价并不相同。现有系统通常只报告准确率或F1，或者只报告可信集覆盖率，缺少在同一损失下比较点答案、裁决集和弃权的统一口径。可信集覆盖率高并不等于决策损失低，集合宽度和弃权成本需要进入评价。

上述三个挑战对应本文的三项设计。证据表示由三层锚点和五类资格判定回应；时间演化由公开揭示日约束的Beta可靠性回应；决策口径由全动作损失感知读出回应。三项设计共同构成FinPlanRAG框架，其核心是AnchorEvolve锚点演化模块。

{\sloppy
还需要说明基准缺口。公开基准中，DocFinQA、FinQA、TAT-QA和FinanceBench的答案已经存在于文档中，FinDVer和Fin-Fact的断言已经给定、标签是支持、反驳或证据不足。没有一个公开基准同时包含断言、截止时点、结果空间、时点证据和已实现结果五元组。本文因此构建收入指引断言基准，并明确说明公开基准的适配边界。公开基准可以检验证据检索、阅读器和断言资格判定等组件，但不能直接检验前瞻性结果预测。把公开基准的数字包装成端到端优势，既不符合任务定义，也会误导读者。\par
}

% ==================== 第3章 ====================
\chapter{问题定义与研究设计}
\label{ch:problem}

\section{前瞻性断言契约}

本文把一条前瞻性断言表示为一个契约
\begin{equation}
q=(e,p,m,\mathcal H,c),
\label{eq:contract}
\end{equation}
其中$e$是公司，$p$是目标期间，$m$是指标，$\mathcal H$是有限结果空间，$c$是截止时点。契约的作用是把“回答什么”和“何时回答”写成显式对象。不同的公司、期间和指标可以共享同一套算法，但契约本身不隐藏。

结果空间$\mathcal H=\{h_1,\dots,h_K\}$是互斥且穷尽的结果标签集合。对于数值指标，标签由阈值定义。例如，营业收入同比增速可以划分为“低于140\%”“140\%至170\%”“高于170\%”，归母净利润同比增速可以划分为“低于190\%”“190\%至220\%”“高于220\%”，综合毛利率可以划分为“低于45\%”“45\%至47\%”“高于47\%”。阈值可以来自上一年度实际值、已披露的季度数据、分析师预测或用户给定的风险边界。阈值的选择会影响结论，因此必须写入契约并在审计报告中展示。

对于披露日期，标签由日历窗口定义。例如，2026年年度报告的披露日期可以划分为“不晚于2027-03-31”“2027-04-01至2027-04-15”“晚于2027-04-15”。窗口的边界来自交易所规则、公司历史披露习惯和预约披露安排。日期断言的结果空间不是数值区间，而是日历区间；第\ref{ch:algorithm}章给出日期标签的投影和估计方法。

截止时点$c$是契约的时间边界。可用证据集合为
\begin{equation}
V_c=\{e:a(e)\le c,\ e\text{在}c\text{时未被撤回}\},
\label{eq:cutoff}
\end{equation}
其中$a(e)$是证据的公开可用时间。对于申报文件，$a(e)$取文件提交日；对于公告，取公告日；对于网页，取页面发布日。发布日期未知的证据不能作为硬锚点，其可靠性上限由配置给出。晚于$c$的证据在进入锚点抽取前被剔除，并写入审计报告的剔除列表。

契约还包含一个隐含假设：目标期间的结果在$c$时尚未公开。如果结果已经公开，问题就不再是前瞻性断言，而应当由确定性检索回答。本文的系统在契约构造阶段检查这一点；如果证据中出现目标期间的实际结果，系统应提示用户该断言已经实现，而不是继续做前瞻性判断。

\section{结果空间与损失}

给定契约$q$，系统维护结果空间上的信念$b\in\Delta(\mathcal H)$，其中$b(h)$是对结果$h$的概率。系统的动作空间为
\begin{equation}
\mathcal A=\mathcal H\cup\{R\subseteq\mathcal H:|R|\ge2\}\cup\{\text{abstain}\}.
\label{eq:action}
\end{equation}
动作空间包含三类动作：点答案$h\in\mathcal H$，非单元素裁决集$R$，以及弃权。单元素裁决集与点答案在语义上重复，因此不进入动作空间。这个限制使三类动作的含义保持清晰：点答案是明确判断，裁决集是保留多个可能结果，弃权是不给出结果。

损失函数分为三部分。点答案的损失为
\begin{equation}
L_{\text{ans}}(h,y)=
\begin{cases}
0, & h=y,\\
\ell_{\text{over}}, & h=h_K,\ y=h_1,\\
\ell_{\text{under}}, & h=h_1,\ y=h_K,\\
\ell_{\text{wrong}}, & \text{其他},
\end{cases}
\label{eq:ansloss}
\end{equation}
其中$h_1$和$h_K$是结果空间的下端和上端。过度乐观的代价高于过度悲观，因此$\ell_{\text{over}}>\ell_{\text{under}}$。裁决集的损失为
\begin{equation}
L_{\text{set}}(R,y)=\indic[y\notin R]\,\ell_{\text{miss}}
+\ell_{\text{width}}\max(0,|R|-1),
\label{eq:setloss}
\end{equation}
其中第一项惩罚漏掉真实结果，第二项惩罚集合宽度。弃权的损失为常数$\ell_{\text{abstain}}$。默认参数为$\ell_{\text{over}}=12$、$\ell_{\text{under}}=8$、$\ell_{\text{wrong}}=10$、$\ell_{\text{miss}}=8$、$\ell_{\text{width}}=2$、$\ell_{\text{abstain}}=5$。这些参数对应金融分析中对过度乐观、漏报和集合宽度的相对定价，不是从测试集调出来的。不同机构可以根据自身风险偏好替换参数，而不改变框架结构。

给定信念$b$，动作的期望损失为
\begin{equation}
\E_{y\sim b}[L(a,y)]=\sum_{h\in\mathcal H}b(h)L(a,h).
\label{eq:expectedloss}
\end{equation}
系统选择期望损失最小的动作
\begin{equation}
a^\star=\arg\min_{a\in\mathcal A}\E_{y\sim b}[L(a,y)].
\label{eq:argmin}
\end{equation}
这个读出是损失感知的，因为它把不同动作放在同一个损失下比较，而不是先固定一个阈值再决定是否弃权。

\section{为什么单一答案不足}

单一答案在前瞻性断言上有三个问题。

第一，证据通常只支持部分结果。截止时点的证据可能表明结果不太可能落在低端，却不足以区分中端和高端。此时点答案会丢掉中端和高端之间的不确定性，裁决集会保留两者。损失函数通过集合宽度和漏报惩罚来定价这种保留。

第二，错误成本不对称。金融分析对过度乐观的代价通常高于过度悲观，因为过度乐观可能导致高估收益、低估风险。点答案的损失函数已经把这种不对称写进去；裁决集和弃权的损失进一步允许系统在不确定性高时选择更保守的动作。

第三，可信集不等于决策集。$1-\alpha$可信集是最小覆盖$1-\alpha$概率质量的集合，它的目标是覆盖率；损失最优裁决集是最小化期望损失的集合，它的目标是决策风险。两者可以在同一个后验上给出不同的集合。把可信集直接当作决策集，会把覆盖率当成决策目标，忽略集合宽度和弃权成本。本文同时报告两者，但把损失最优裁决集作为决策口径。

\section{研究问题与假设}

本文的研究问题可以写成四个可检验的假设。

\textbf{H1（锚点的语义增量）}：带来源的文本锚点能够为纯LLM分支提供可测量的语义信号。检验口径为纯LLM锚点后验与无锚点或少锚点条件的准确率、平衡准确率和宏F1。若带锚点条件在多数基座上不高于无锚点条件，则H1不成立。

\textbf{H2（资格判定与时间演化的校准作用）}：五类资格判定和按揭示日更新的Beta可靠性能够降低决策损失，并改善后验的稳定性。检验口径为静态锚点与演化锚点的损失差、不同锚点类型的消融、以及无日期来源和未通过引文核验的锚点比例。若演化条件在多数基座上不优于静态条件，则H2不成立。

\textbf{H3（全动作读出的决策改进）}：在同一个后验和同一个损失下，全动作损失感知读出严格优于受限可信集门。检验口径为三个基座上的配对损失差和95\%置信区间。若任一基座上的置信区间包含零或方向相反，则H3不成立。

\textbf{H4（与强结构化基线的竞争性）}：包含锚点状态的系统在决策损失上与强随机森林诊断相当，其价值主要体现在可审计的语义状态，而不是显著的点预测优势。检验口径为最终系统与演化先验随机森林的损失差及置信区间。若最终系统显著劣于随机森林，则H4不成立；若最终系统显著优于随机森林，则本文可以进一步主张锚点层的点预测增量。

四个假设对应四个层次：H1检验锚点是否提供语义信号，H2检验资格判定和时间演化是否改善校准，H3检验读出是否改善决策，H4检验系统与强结构化基线的相对位置。这样的假设设计使本文能够区分“锚点有用”“演化有用”“读出有用”和“整体系统更强”四个不同主张，而不是把它们混为一个端到端结论。

\section{总体框架}

FinPlanRAG由六个模块组成：证据构造、锚点抽取、资格判定、时间演化、信念更新和全动作读出，其中锚点抽取、资格判定、时间演化和信念更新构成AnchorEvolve模块。图\ref{fig:framework}给出总体结构。

\begin{figure}[htbp]
\centering
\begin{tikzpicture}[
  font=\small,
  >=Latex,
  box/.style={draw, rounded corners=2pt, align=center, text width=3.2cm, minimum height=1.0cm, inner sep=2pt},
  small/.style={draw, rounded corners=2pt, align=center, text width=2.8cm, minimum height=0.9cm, inner sep=2pt},
  arrow/.style={->, thick, draw=black!70}
]
\node[box] (q) at (0,0) {断言契约\\$q=(e,p,m,\mathcal H,c)$};
\node[box] (ev) at (0,-1.7) {时点证据\\公告、申报文件、网页};
\node[box, fill=gray!8] (ax) at (0,-3.4) {锚点抽取\\L0/L1/L2};
\node[small] (h) at (-3.9,-5.2) {硬锚点\\约束可行集};
\node[small] (s) at (0,-5.2) {软锚点\\似然更新};
\node[small] (f) at (3.9,-5.2) {空泛/过期/对抗\\过滤};
\node[box] (b) at (0,-6.9) {结果信念 $b$};
\node[box] (st) at (-3.9,-8.6) {时间演化\\Beta可靠性/先验};
\node[box] (ba) at (0,-8.6) {锚点后验 $b_A$};
\node[box, fill=gray!8] (br) at (3.9,-8.6) {结构化RF后验 $b_R$};
\node[box] (bf) at (0,-10.3) {校准融合 $b_F$};
\node[box] (out) at (0,-12.0) {全动作损失感知读出\\点答案/裁决集/弃权};
\node[box] (au) at (5.2,-12.0) {审计文件\\证据/锚点/损失};
\draw[arrow] (q) -- (ev);
\draw[arrow] (ev) -- (ax);
\draw[arrow] (ax) -- (h);
\draw[arrow] (ax) -- (s);
\draw[arrow] (ax) -- (f);
\draw[arrow] (h) -- (b);
\draw[arrow] (s) -- (b);
\draw[arrow] (b) -- (ba);
\draw[arrow] (st) -- (ba);
\draw[arrow] (ba) -- (bf);
\draw[arrow] (br) -- (bf);
\draw[arrow] (bf) -- (out);
\draw[arrow] (out) -- (au);
\end{tikzpicture}
\caption{FinPlanRAG的总体框架。断言契约把公司、期间、指标、结果空间和截止时点写成显式对象；时点证据经过锚点抽取和五类资格判定后进入信念更新；时间演化状态只使用已公开揭示的结果；锚点后验与结构化随机森林后验融合；全动作损失感知读出选择点答案、裁决集或弃权；全部中间状态写入审计文件。}
\label{fig:framework}
\end{figure}

证据构造模块从公告、申报文件和网页中获取材料，按截止时点过滤，分配来源编号。锚点抽取模块调用LLM，要求每条锚点引用来源编号、给出逐字引文、发布日期、隐含结果标签和校准可靠性。资格判定模块用确定性规则检查空泛、过期、对抗、硬、软五类状态。时间演化模块按锚点类型维护Beta可靠性，只用已公开揭示的结果更新。信念更新模块把软锚点转换为似然，把硬锚点转换为约束，得到锚点后验。全动作读出模块把锚点后验与随机森林后验融合，在点答案、裁决集和弃权上选择期望损失最小的动作。

六个模块之间通过结构化对象连接。证据是带来源编号的文本对象，锚点是带类型、来源、日期、可靠性、隐含标签的JSON对象，信念是结果空间上的概率向量，动作是标签、集合或弃权。每个对象都可以单独检查，模块之间不共享隐状态。这个设计使系统可以在不重新训练端到端模型的情况下替换证据源、锚点类型或校准后端。

\section{时点约束与评价协议}

时点约束是本文的核心设计约束。设断言$q$的截止时点为$c$，可用证据集合为$V_c$。系统在$c$时只能使用$V_c$中的证据，不能使用$c$之后公布的季度结果、年报、分析师更新或价格数据。这个约束适用于证据构造、锚点抽取、可靠性更新、先验估计、特征构造、校准和权重选择。

具体而言，本文采用以下规则。第一，证据的发布日期必须不晚于$c$；发布日期未知的证据不能作为硬锚点，可靠性上限为0.55。第二，锚点类型的Beta可靠性只使用揭示日早于$c$的已揭示结果；揭示日取财报事件日期，缺失时回退到10-Q日期。第三，结果先验只使用已揭示结果的平滑计数；没有已揭示结果时退化为均匀先验。第四，随机森林的训练样本只来自其他公司折，且结果日期早于测试年份1月1日；测试公司不出现在训练集中。第五，温度缩放和融合权重只在验证集上拟合或选择，测试集不参与。第六，所有LLM请求按请求负载的哈希缓存，相同输入产生相同输出；搜索结果不缓存，因为网页会变化，审计文件记录当次使用的证据。

评价协议与假设对应。H1和H2在227条XBRL核验测试断言上评估，报告准确率、平衡准确率、宏F1和决策损失。H3在三个基座上比较受限可信集门和全动作读出，报告配对损失差和95\%置信区间。H4比较最终系统与演化先验随机森林，报告损失差、置信区间和McNemar检验。所有置信区间使用10,000次断言级配对自助法；准确率差异使用McNemar检验。由于测试集只有227条断言，本文不把单个点估计的差异解释为显著优势，只有在置信区间排除零时才主张显著改进。

% ==================== 第4章 ====================
\chapter{FinPlanRAG框架}
\label{ch:framework}

\section{设计原则}

FinPlanRAG的设计遵循六条原则，其核心模块AnchorEvolve承担证据状态化、资格判定、时间演化和信念更新。这些原则不是附加的功能列表，而是对第\ref{ch:problem}章所述研究缺口的直接回应。

第一，契约先行。系统在检索证据之前先确定公司、目标期间、指标、结果空间和截止时点。契约把“回答什么”写成显式对象，使不同方法可以在同一口径下比较，也使用户可以在运行前检查阈值是否合理。契约不是系统的内部细节，而是审计报告的第一部分。

第二，来源落地。每条锚点必须引用一个具体的证据来源，并给出逐字引文。系统不接受没有来源的判断，也不接受只在理由中描述证据而不给出引文的锚点。来源落地使锚点可以被逐条复核，也使错误的锚点可以被定位到具体的网页或段落。

第三，资格先于权重。锚点在进入信念更新之前先经过资格判定。空泛、过期、对抗三类锚点不进入更新；硬锚点和软锚点分别处理。这个顺序很重要：如果先把所有锚点加权再过滤，被拒绝的锚点仍然会影响后验；如果先过滤再加权，每个被采纳的锚点都有明确的资格理由。

第四，揭示后演化。锚点类型的可靠性只在结果公开揭示之后更新，且只使用揭示日早于当前截止时点的结果。这个约束使可靠性演化在回测和部署中具有相同的含义，避免把未来结果泄漏到时点判断。

第五，损失感知。输出不是固定的点答案，而是在点答案、非单元素裁决集和弃权三个动作上最小化期望损失。损失函数把过度乐观、漏报、集合宽度和弃权成本写成显式参数，使不同风险偏好可以在同一框架下表达。

第六，审计优先。系统把证据、锚点、资格、可靠性、信念轨迹、动作损失和日期估计写入结构化文件。审计不是事后附加的日志，而是系统输出的一部分。没有审计文件的判断不能被本文视为可复核的判断。

\section{三层锚点状态}

FinPlanRAG的核心状态是AnchorEvolve的三层锚点。三层分别是稳定关系层（L0）、行业模式层（L1）和实例层（L2）。这里的锚点不是知识图谱中的实体或关系，也不是行为金融中的anchoring bias；L0、L1和L2分别对应关系约束、行业模式绑定和带来源的证据实例，锚定的是证据对结果的约束关系，而不是一个实体节点。

L0是稳定关系本体。它定义一组可复用的关系类型，包括界、先后、前置时间、持续时间、守恒、比率、单位换算和作用域。L0不随行业变化，为锚点提供语义骨架。例如，“当前季度值落在指引区间内”是一个界关系，“新产品在指引季度之后投产”是一个先后关系，“产能爬坡需要两个季度”是一个前置时间关系，“分部收入之和等于总收入”是一个守恒关系，“毛利率等于毛利润除以营业收入”是一个比率关系。L0的作用是限制锚点可以表达的关系种类，使不同行业的锚点共享同一套语义。

L1是行业模式。每个锚点类型是一个四元组
\begin{equation}
\tau=(\text{名称},\ \text{L0关系},\ \text{字段},\ \text{生效轮次}),
\label{eq:l1}
\end{equation}
其中字段给出该类型需要从证据中读取的信息，生效轮次给出该类型可以在哪一轮断言中使用。收入指引的初始模式包括指引宽度、当前季度增长、公司历史、行业历史、文本信号、产能与需求、分析师预测、披露历史和监管规则。L1把L0关系绑定到具体字段和生效规则。例如，“指引宽度”绑定到界关系，字段是上限与下限之差除以中点；“当前季度增长”绑定到比率关系，字段是当前季度值相对指引中点的比例；“披露历史”绑定到先后关系，字段是公司过去若干年年报的公告日期。

L2是锚点实例。每条锚点是一个元组
\begin{equation}
a=(id,\tau,v,u,s,src,d,r,\text{hard},\hat h,\text{note}),
\label{eq:l2}
\end{equation}
其中$id$是锚点编号，$\tau$是L1类型，$v$是取值或区间，$u$是单位，$s$是作用域，$src$是来源编号，$d$是来源日期，$r$是可靠性，$\text{hard}$是硬锚点标记，$\hat h$是隐含结果标签，$\text{note}$是理由。L2是LLM抽取器的直接输出，也是资格判定和信念更新的输入。每条L2锚点都记录来源和时间，因此可以被逐条检查。

三层的关系可以概括为：L0定义可以说哪些关系，L1定义某类断言需要哪些字段，L2记录某条证据的具体取值。三层不是三个独立训练的模块，而是同一状态的不同抽象层次。L1类型用于可靠性演化，L0关系用于语义分类，L2实例进入信念更新。这样的分层使系统可以在不改变L0的情况下增加新的L1类型，也可以在不改变L1的情况下替换L2抽取器。

\section{锚点类型}

表\ref{tab:anchortypes}给出本文使用的主要锚点类型。类型的选择遵循两个标准：一是该类型在金融文本中有稳定的语言表达，二是该类型能够对应到结果空间中的一个或多个标签。类型不是越多越好；过多的类型会增加抽取器的负担，也会增加可靠性后验的稀疏性。

\begin{table}[H]
\centering
\zihao{5}
\caption{FinPlanRAG的主要锚点类型。}
\label{tab:anchortypes}
\begin{tabular}{p{2.6cm}p{2.4cm}p{7.2cm}}
\toprule
类型 & L0关系 & 作用 \\
\midrule
指引宽度 & 界 & 指引区间越宽，结果落在区间内的先验越高；区间本身不是结果落在区间内的证据。\\
当前季度增长 & 比率 & 当前季度实际值相对指引中点的比例，用于判断全年结果的起点。\\
公司历史 & 先后 & 公司过去已揭示结果相对指引的偏差，用于估计公司层面的兑现倾向。\\
行业历史 & 先后 & 同行业其他公司已揭示结果的分布，用于估计行业层面的共同冲击。\\
文本信号 & 语义 & 管理层评论、需求描述、订单和产能信息，用于补充结构化字段。\\
产能与需求 & 前置时间 & 产能投放、利用率和订单覆盖，用于判断收入增长的可持续性。\\
分析师预测 & 界 & 卖方机构对目标期间的预测，作为软锚点而不是事实。\\
披露历史 & 先后 & 公司过去年度报告的公告日期，用于披露日期断言。\\
监管规则 & 界 & 交易所的披露截止日，作为上界而不是具体日期证据。\\
\bottomrule
\end{tabular}
\end{table}

锚点类型与结果标签的关系不是一一对应。一条“当前季度增长”锚点可能支持“高于”标签，也可能只支持“不低于”某个下界；一条“指引宽度”锚点通常不直接支持任何标签，而是影响区间内的先验。资格判定和信念更新需要分别处理这些情况。锚点的$\hat h$字段记录抽取器认为该锚点隐含的标签；如果锚点不提供标签信息，$\hat h$为空，锚点被判为空泛。

\section{证据来源与作用域}

FinPlanRAG使用三类证据来源。

第一类是公司申报文件和公告。美国公司使用8-K、10-Q、10-K和XBRL数据；中国A股公司使用巨潮资讯网的公告记录、年报、季报和临时公告。申报文件提供结构化事实和原文段落，公告记录提供披露日期和文件链接。官方来源的优先级高于媒体转载，因为官方来源的日期和文本更可靠。

第二类是财经媒体和研究报告。媒体文章可以提供管理层评论、行业需求、产能和价格信息，但可能存在转载、改编和日期错误。系统要求媒体来源给出逐字引文，并对无日期来源降权。研究报告中分析师预测作为软锚点，不作为事实。

第三类是交易所规则和监管文件。这类来源提供披露截止日、披露格式和合规要求。监管规则通常只给出上界，不能直接支持某个具体日期窗口；系统把泛泛的监管截止日判为空泛锚点，除非来源明确给出公司的预约披露日期。

锚点的作用域分为公司、行业、市场和监管四类。作用域的作用是防止跨作用域误用。例如，同行业其他公司的营收指引结果可以作为行业锚点，但不能直接当作目标公司的结果；交易所的披露截止日可以作为监管锚点，但不能直接推断公司会在截止日当天披露。资格判定检查作用域是否与断言匹配，不匹配的锚点被判为对抗锚点。

\section{与既有框架的区别}

FinPlanRAG与断言验证、RAG和选择性预测的区别可以概括为三点。

第一，状态不同。断言验证的输出是一个离散裁决，RAG的输出是一个生成答案，选择性预测的输出是一个标签或弃权。FinPlanRAG的输出是一个可演化的锚点状态，包含证据、资格、可靠性和信念。这个状态可以被逐条检查，也可以在结果揭示后更新。

第二，时间不同。断言验证和RAG通常在一次调用中完成，不区分证据的可用时间和结果的揭示时间。FinPlanRAG把截止时点写进契约，把揭示日写进可靠性更新规则。这个区分使系统可以在回测中复现部署时的信息条件。

第三，决策不同。断言验证和RAG通常报告准确率或F1，选择性预测报告覆盖率或弃权率。FinPlanRAG在同一个损失下比较点答案、裁决集和弃权，报告决策损失、覆盖率和集合大小。这个口径使系统可以在不确定性高时选择集合或弃权，而不是被迫输出一个点答案。

% ==================== 第5章 ====================
\chapter{算法设计}
\label{ch:algorithm}

本章给出FinPlanRAG中AnchorEvolve模块的完整算法设计。本章第1节到第7节分别介绍证据构造、锚点抽取、资格判定、时间演化、信念更新、全动作读出和披露日期断言，最后给出算法伪代码、复杂度和审计格式。

\section{证据构造}

证据构造的目标是为每条断言提供一组带日期、带来源编号、可引用的材料。构造过程分为查询生成、检索、去重、日期过滤和编号五步。

查询生成从契约出发。给定公司$e$、目标期间$p$和指标$m$，系统生成四类查询：公司加期间加指标，公司加最近季度加指标，公司加上一年度加指标，以及公司加“预测”“研报”“公告”等限定词。对于披露日期断言，系统额外生成公司加“年度报告披露日期”、公司加“往年披露时间”、公司加“预约披露”等查询。查询数量由配置限制，默认每条断言最多四条。查询中不放入断言的目标答案，避免把判断泄漏到检索。

检索使用多个网页搜索提供者，并按轮换顺序调用。某个提供者返回空结果或安全验证页时，系统自动切换到下一个提供者。这个设计是为了避免单一搜索引擎的限流让整次运行静默失败。检索结果包含标题、链接、摘要、发布日期和提供者。对于披露日期断言，系统还查询巨潮资讯网的公告接口，直接获取公司历年年报的公告日期和PDF链接。官方公告记录作为高优先级证据，排在网页结果之前。

去重按链接和标题进行。同一篇文章的不同转载只保留一条；同一链接的不同摘要只保留一条。日期过滤检查证据的发布日期是否晚于截止时点。晚于截止时点的证据被写入剔除列表，不进入锚点抽取。发布日期未知的证据保留，但标记为无日期，并在资格判定中降权。

编号使用$W1,W2,\dots$，按证据得分排序。得分由日期是否已知、摘要是否包含数字、标题或摘要是否包含公司名、文本长度和查询顺序决定。官方公告记录得分最高。编号的作用是让锚点可以引用具体来源，也让审计报告可以把锚点映射回证据。

对于申报文件，系统在证据构造阶段还做段落切分。每个段落带稳定编号$P1,P2,\dots$，结构化证据带$E1,E2,\dots$。抽取器可以引用$P\#$或$E\#$，资格判定检查引用的段落是否在时点摘录中。这个设计使锚点可以定位到段落，而不是只定位到文件。

\section{锚点抽取}

锚点抽取是一个LLM调用。输入包括断言契约、结果标签及其定义、参考口径、时点证据、锚点类型和关系类型。输出是一个JSON对象，包含锚点列表、标签分数、日期估计（仅日期断言）、置信度和缺失说明。

提示词要求锚点遵守以下规则。第一，每条锚点必须引用一个证据编号，不能引用不存在的来源。第二，每条锚点必须给出逐字引文，引文必须在来源文本中出现。第三，锚点的发布日期不得晚于截止时点；来源日期未知时，日期设为空，可靠性不得高于0.55。第四，可靠性是“隐含标签正确的校准概率”，不是推理置信度。第五，硬锚点必须是显式约束，能够排除至少一个标签；管理层评论和分析师预测通常是软锚点。第六，证据不足时在缺失字段中说明，不编造事实。第七，标签分数必须使用契约中的标签，且和为1。

抽取器的输出是一个结构化对象，而不是自由文本。这个约束使后续的资格判定可以按字段检查，而不是依赖另一个LLM来判断前一个LLM的输出是否合理。解析失败的响应被记录为错误，不进入信念更新；系统不使用金标结果作为回退。这个设计是为了避免在评估中把解析失败悄悄地替换成正确答案。

锚点抽取的成本主要来自输入和输出长度。单条断言的输入通常在数千token量级，输出在数百token量级。系统按请求负载的哈希缓存响应，相同输入产生相同输出。缓存不包含API密钥，也不写入审计文件的公开部分。

\section{五类资格判定}

资格判定是FinPlanRAG与普通RAG的关键区别。普通RAG把检索到的片段直接交给生成器；FinPlanRAG在锚点进入信念更新之前，先用确定性规则检查锚点的资格。表\ref{tab:qualification}给出五类判定。

\begin{table}[H]
\centering
\zihao{5}
\caption{五类锚点资格判定。}
\label{tab:qualification}
\begin{tabular}{p{2.1cm}p{7.0cm}p{4.2cm}}
\toprule
类型 & 判定条件 & 处理 \\
\midrule
空泛 & 没有隐含标签，或锚点不提供有效标签信息 & 不进入更新 \\
过期 & 锚点日期或来源日期晚于截止时点 & 不进入更新 \\
对抗 & 来源编号不在证据集合中，或作用域不匹配 & 不进入更新 \\
硬 & 显式约束、引文通过、来源日期已知、可靠性不低于阈值 & 约束可行集 \\
软 & 其余有效锚点 & 似然更新 \\
\bottomrule
\end{tabular}
\end{table}

资格判定的具体规则如下。

第一，来源检查。锚点的来源编号必须出现在本次证据集合中。来源不存在，判为对抗。这个检查防止LLM引用一个看起来合理但实际不存在的来源。

第二，日期检查。锚点日期或来源日期晚于截止时点，判为过期。如果锚点日期与来源日期不一致，系统以来源日期为准检查过期，并对可靠性乘以0.9。日期不一致本身不是拒绝理由，但它是一个时间一致性疑点。

第三，引文检查。锚点的引文必须在来源文本中逐字出现。系统对引文和来源文本做空白规范化，然后做子串匹配。引文为空、长度过短或匹配失败，判为引文未通过。未通过的锚点可靠性乘以0.7，且不能成为硬锚点。这个检查是确定性的，不依赖另一个LLM。

第四，作用域检查。锚点的作用域必须是公司、行业、市场或监管之一。作用域不匹配，判为对抗。例如，把同行业其他公司的结果当作目标公司的直接结果，作用域不匹配。

第五，无日期检查。锚点日期和来源日期都未知时，锚点保留为软锚点，但可靠性上限为0.55，且不能成为硬锚点。这个规则承认无日期证据可能有用，但不允许它主导判断。

第六，参考值检查。对于增速类断言，只描述上一年度参考值的锚点判为空泛。上一年度实际值是契约的参考口径，不是目标期间的结果证据。这个检查防止抽取器把基数当作结果，从而把信念拉向低增长。

第七，监管规则检查。对于披露日期断言，泛泛的法定截止日判为空泛。截止日是一个上界，不能区分三个日期窗口；只有公司预约披露日期或历史披露记录才能支持具体窗口。

第八，日期标签投影。对于披露日期断言，如果锚点提到的日期与隐含标签不一致，系统把日期中的月和日投影到目标年份，再匹配日期窗口。例如，历史年报日期2025-04-21投影到2027-04-21，落在“晚于2027-04-15”窗口。投影只在一个锚点提到的所有日期指向同一窗口时生效；如果锚点同时提到多个窗口的日期，系统不覆盖模型标签，只记录不一致。这个规则修正了抽取器在跨年日期上的常见错误。

第九，时间一致性检查。对于披露日期断言，如果引文中提到的日期晚于来源发布日期，且不晚于截止时点，系统标记时间一致性疑点，可靠性乘以0.5，且不能成为硬锚点。这个检查处理动态网页或转载文章在原始发布日期之后被更新，导致引文包含当时不可能知道的信息的情况。

资格判定的输出是每条锚点的状态、资格类型、有效可靠性、引文是否通过、时间一致性标记和理由。被拒绝的锚点保留在审计文件中，理由逐条列出；被采纳的锚点进入信念更新。这个设计使审计者可以看到系统拒绝了什么，以及为什么拒绝。

\section{时间演化}

时间演化的目标是让锚点类型的可靠性随已揭示结果更新，同时不把未来信息泄漏到时点判断。设断言$q$的截止时点为$c$，锚点类型为$\tau$。截至$c$已揭示的、类型为$\tau$的锚点中，隐含标签与结果一致的数量为$m_\tau$，总数为$n_\tau$。类型$\tau$的Beta可靠性为
\begin{equation}
r_\tau=\frac{m_\tau+\alpha}{n_\tau+\alpha+\beta},
\label{eq:beta}
\end{equation}
其中$\alpha,\beta$是Beta先验参数。默认取$\alpha=\beta=1$，即均匀先验。当$n_\tau=0$时，$r_\tau=0.5$。

抽取器给出的可靠性是LLM对单条锚点的主观判断，不是经过验证的频率。系统把演化可靠性与LLM可靠性做收缩：
\begin{equation}
\rho_\tau=\lambda r_\tau+(1-\lambda)r_{\mathrm{LLM}},
\label{eq:shrink}
\end{equation}
其中$\lambda$是收缩系数，默认0.5。$\rho_\tau$作为软锚点更新的指数。这个收缩的作用是：当历史样本少时，系统更多依赖LLM判断；当历史样本多时，系统更多依赖已揭示结果的频率。

先验也按揭示日更新。设截至$c$已揭示的结果中，落在标签$h$的数量为$n_h$，总数为$n$，则先验为
\begin{equation}
b_0(h)=\frac{n_h+1}{n+|\mathcal H|}.
\label{eq:prior}
\end{equation}
当$n=0$时，先验为均匀分布。这个先验只使用已揭示结果，不使用目标期间的结果。

演化状态按轮次更新。第$r$轮的结果在第$r$轮全部断言裁决完成后才更新状态，并从第$r+1$轮开始使用。新发现的锚点类型的生效轮次为$r+1$，因此触发发现的断言不能使用该类型。这个规则保证自演化是时点正确的。本文的主实验报告单轮演化；多轮循环作为扩展接口保留，第\ref{ch:conclusion}章讨论其局限。

\section{锚点类型自演化}

固定锚点类型库依赖人工设计，难以覆盖不同指标和不同市场的证据模式。本文在Beta可靠度演化之外增加一个受验证集约束的类型自演化循环。它不是让模型自由改写标签或测试集，而是让模型在验证误差上提出新的锚点类型，只有通过验证集损失检验的类型才进入类型库。

设第$t$轮的锚点类型库为$\mathcal T_t$，验证集上由当前类型库得到的决策损失为$L_{\mathrm{val}}(\mathcal T_t)$。循环分四步。

第一步，误差挖掘。把验证集按公司不重叠地分成挖掘集和选择集。在挖掘集上，用当前类型库计算每条断言的预测和损失，选出预测错误或损失较高的断言。挖掘集只用于提出候选类型，不参与候选接受。

第二步，类型提出。把错误断言、对应证据、当前锚点和真实结果交给一个LLM，要求它提出至多$k$个新的锚点类型。每个候选类型是一个四元组
\begin{equation}
\tau'=(\text{名称},\ \text{L0关系},\ \text{字段},\ \text{适用条件}),
\label{eq:candidatetype}
\end{equation}
并附带一句证据示例。提出器被要求只描述可复用的证据模式，不能引用具体测试断言，不能把结果标签本身当作类型。

第三步，候选重抽取。把候选类型加入类型库，重写锚点抽取提示词中的类型说明，在选择集上重新抽取锚点并计算后验。类型库扩展为$\mathcal T_t\cup\mathcal T'$。

第四步，验证门控。只有当选择集上的损失下降超过阈值$\delta$时，才接受候选类型：
\begin{equation}
L_{\mathrm{val}}(\mathcal T_t)-L_{\mathrm{val}}(\mathcal T_t\cup\mathcal T')>\delta.
\label{eq:typeaccept}
\end{equation}
否则拒绝整组候选，保持原类型库。接受后进入下一轮，最多进行$R$轮。测试集只在循环冻结后评估一次，不参与类型提出、阈值选择或轮数选择。

这个循环与提示词演化分开评估。提示词改写、类型扩展和二者组合分别作为候选臂，在同一选择集上比较；只有损失最低且超过阈值的臂被接受。这样可以把“模型提出了新类型”和“新类型真的改善了验证损失”分开，避免把语言上的合理性当作经验收益。

自演化的代价是额外的验证集抽取。每轮需要重抽选择集锚点，但不重抽测试集；类型库冻结后才抽取测试集。由于循环只使用验证集，它与主实验的时点约束一致，也不会把测试结果泄漏到类型库。

类型库扩展是框架的训练步骤，不是独立模块。算法\ref{alg:typeexpand}给出伪代码，其输出$\mathcal T^\star$作为算法\ref{alg:anchorevolve}的类型库输入。主实验中DeepSeek使用通过验证门控的类型库，GLM和Qwen保留基础库；这样，类型库的选择和其他超参数一样，只由验证集决定。

\begin{algorithm}[H]
\caption{类型库的验证集扩展}
\label{alg:typeexpand}
\begin{algorithmic}[1]
\Require 类型库$\mathcal T_0$，验证集$\mathcal D_{\mathrm{val}}$，阈值$\delta$
\Ensure 冻结类型库$\mathcal T^\star$
\State 将$\mathcal D_{\mathrm{val}}$按公司不重叠划分为挖掘集$\mathcal D_{\mathrm{mine}}$和选择集$\mathcal D_{\mathrm{select}}$
\State 在$\mathcal D_{\mathrm{mine}}$上计算当前类型库的损失，选出高损失断言
\State 由LLM提出候选类型$\mathcal T'$，每个类型包含关系、字段和适用条件
\State 在$\mathcal D_{\mathrm{select}}$上重抽锚点，计算$L_{\mathrm{val}}(\mathcal T_0\cup\mathcal T')$
\If{$L_{\mathrm{val}}(\mathcal T_0)-L_{\mathrm{val}}(\mathcal T_0\cup\mathcal T')>\delta$}
\State $\mathcal T^\star\gets\mathcal T_0\cup\mathcal T'$
\Else
\State $\mathcal T^\star\gets\mathcal T_0$
\EndIf
\State \Return $\mathcal T^\star$ \Comment{冻结后用于测试集}
\end{algorithmic}
\end{algorithm}


\section{信念更新}

信念从先验$b_0$开始。对于软锚点$a$，设其隐含标签为$\hat h$，有效可靠性为$\rho$，软似然为
\begin{equation}
L_a(h)=
\begin{cases}
p, & h=\hat h,\\
\dfrac{1-p}{|\mathcal H|-1}, & h\ne\hat h,
\end{cases}
\label{eq:softlikelihood}
\end{equation}
其中$p$是配置项，默认0.7。软锚点的更新为
\begin{equation}
b(h)\propto b(h)\,L_a(h)^{\rho}.
\label{eq:softupdate}
\end{equation}
指数$\rho$控制锚点的影响强度：可靠性越高，更新越强；可靠性越低，更新越弱。更新后对信念做归一化。

对于硬锚点，系统把信念压到隐含标签上，并保留一个小底：
\begin{equation}
b(h)=
\begin{cases}
1-\epsilon, & h=\hat h,\\
\dfrac{\epsilon}{|\mathcal H|-1}, & h\ne\hat h,
\end{cases}
\label{eq:hardupdate}
\end{equation}
其中$\epsilon$默认0.02。小底的作用是避免后续软锚点无法恢复。多个硬锚点的可行集交集、冲突检测和回退到软更新是论文的设计目标，当前实现只处理单硬锚点的点质量更新；第\ref{ch:conclusion}章把这一点列为局限。

锚点后验记为$b_A$。抽取器还输出标签分数$s$，它是边际信号，不是锚点状态的替代。系统把$b_A$与$s$按权重$\gamma$混合：
\begin{equation}
b_A'=\gamma b_A+(1-\gamma)s,
\label{eq:blend}
\end{equation}
其中$\gamma$默认0.5。混合后的后验保留锚点的结构化信息，同时利用抽取器对三类结果的边际判断。

结构化随机森林后验记为$b_R$。输入特征包括指引宽度、当前季度值相对指引中点的比例、指引与当前季度的比值、历史结果计数与频率、前次指引增长与宽度、目标季度、截止年份、行业，以及演化先验特征（同行业历史、公司意外历史、日历、当前季度相对上下界的位置、可靠性标志）。随机森林在开发集与验证集上训练，并在公司不重叠折上做温度缩放；测试公司不出现在训练集中。

锚点后验与随机森林后验融合：
\begin{equation}
b_F=w b_A'+(1-w)b_R,
\label{eq:fusion}
\end{equation}
其中$w$只在验证集上从配置的权重网格中选择，测试集不参与选择。默认网格为$\{0,0.1,0.2,0.3\}$。最终实验中选择的权重在正文中报告。

\section{全动作损失感知读出}

给定融合后验$b_F$，系统在动作空间$\mathcal A$上选择期望损失最小的动作。动作空间包括$|\mathcal H|$个点答案、$\sum_{k=2}^{|\mathcal H|}\binom{|\mathcal H|}{k}$个非单元素裁决集，以及弃权。对于三标签结果空间，动作空间包含3个点答案、4个非单元素裁决集和1个弃权，共8个动作。

点答案的期望损失为
\begin{equation}
\E[L_{\text{ans}}(h)]=\sum_{y\in\mathcal H}b_F(y)L_{\text{ans}}(h,y).
\label{eq:eans}
\end{equation}
裁决集的期望损失为
\begin{equation}
\E[L_{\text{set}}(R)]=\sum_{y\in\mathcal H}b_F(y)L_{\text{set}}(R,y).
\label{eq:eset}
\end{equation}
弃权的期望损失为$\ell_{\text{abstain}}$。系统选择
\begin{equation}
a^\star=\arg\min_{a\in\mathcal A}\E[L(a)].
\label{eq:fullaction}
\end{equation}

全动作读出与受限可信集门的区别在于动作空间。受限门只比较点答案和最高密度$1-\alpha$可信集，可信集的大小由覆盖率目标决定，而不是由损失决定。全动作读出把点答案、裁决集和弃权放在同一个损失下比较，因此可以在不确定性高时选择裁决集或弃权，也可以在置信度高时选择点答案。对于同一个后验和同一个损失，全动作读出的期望损失不会高于受限门，因为受限门的动作是全动作空间的子集。这个性质是H3的形式基础。

单元素裁决集不进入动作空间。原因是单元素裁决集与点答案在语义上重复，但在损失上不同：如果漏报惩罚低于错误分类惩罚，单元素裁决集会系统性地优于点答案，导致系统把点答案都标成裁决集。这个现象不是决策改进，而是动作标签的重复计数。本文把裁决集限制为两个及以上标签，使三类动作的含义保持清晰。

可信集$R_\alpha$仍然作为覆盖口径报告：
\begin{equation}
R_\alpha=\arg\min_{R\subseteq\mathcal H}\{|R|:\sum_{h\in R}b_F(h)\ge1-\alpha\}.
\label{eq:credible}
\end{equation}
它不参与动作选择，但用于审计和覆盖率报告。损失最优裁决集$R^\star$是动作选择的结果，可以在$R_\alpha$更宽时选择更窄的集合，也可以在不确定性高时选择弃权。两者在审计报告中分开列出，避免把覆盖率误读为决策质量。

\section{披露日期断言}

披露日期断言的算法与数值指标断言共享同一套锚点状态和读出，但结果空间和日期估计不同。

结果空间由日历窗口定义。给定目标期间$p$，系统根据交易所规则和公司历史披露习惯构造三个窗口。例如，2026年年度报告的窗口为“不晚于2027-03-31”“2027-04-01至2027-04-15”“晚于2027-04-15”。窗口必须互斥且覆盖法定披露截止日之前的全部日期。

锚点抽取器除了输出锚点和标签分数，还输出一个日期估计，包括点估计、区间下界、区间上界和依据。日期估计是辅助信息，不直接决定动作；动作仍由标签后验和损失读出决定。系统把日期估计与标签后验做一致性检查：如果点估计落在最高信念窗口内，系统保留点估计；如果点估计落在窗口外，系统用最高信念窗口的中点作为回退点估计。区间下界和上界优先使用模型给出的值，但限制在合理范围内。

日期标签投影是披露日期断言的关键修正。历史年报日期、预约披露日期和监管截止日往往跨年，抽取器容易把历史日期直接映射到当前窗口。系统把日期中的月和日投影到目标年份，再匹配窗口。投影规则只在一个锚点提到的所有日期指向同一窗口时生效，避免把包含多个日期的锚点强行归类。

披露日期断言的评价不使用点日期误差作为主指标，因为目标日期尚未实现，且公告日期与实际披露时点可能相差一天。本文报告动作、裁决集、点估计和区间，并在案例研究中展示审计轨迹。日期断言的定量评价需要等到目标日期实现之后，第\ref{ch:conclusion}章把这一点列为局限。

\section{算法与复杂度}

算法\ref{alg:anchorevolve}给出单条断言的完整流程。

\begin{algorithm}[H]
\caption{AnchorEvolve单条断言的证据构造、锚点演化与全动作裁决}
\label{alg:anchorevolve}
\begin{algorithmic}[1]
\Require 断言契约$q=(e,p,m,\mathcal H,c)$，演化状态$S_c$，类型库$\mathcal T^\star$，随机森林后验$b_R$
\Ensure 动作$a^\star$，点答案$h^\star$，裁决集$R^\star$，后验$b_F$
\State $V_c\gets\textsc{CollectEvidence}(q,c)$ \Comment{按截止时点过滤，分配$W\#$}
\State $A\gets\textsc{ExtractAnchors}(q,V_c,S_c,\mathcal T^\star)$ \Comment{LLM，要求来源与逐字引文}
\State $A_q\gets\textsc{Qualify}(A,q,c)$ \Comment{空泛/过期/对抗/硬/软}
\State $b_A\gets\textsc{BeliefUpdate}(S_c,A_q)$
\State $b_A'\gets\gamma b_A+(1-\gamma)\textsc{LabelScores}(A)$
\State $b_F\gets w b_A'+(1-w)b_R$
\State $(a^\star,h^\star,R^\star)\gets\textsc{FullActionReadout}(b_F)$
\State $S_{c^+}\gets\textsc{RevealUpdate}(S_c,c)$ \Comment{仅用揭示日早于下一截止时点的结果}
\State \Return $(a^\star,h^\star,R^\star,b_F)$
\end{algorithmic}
\end{algorithm}

算法的时间成本由三部分组成。第一是证据构造，包括网页检索、页面抓取和官方公告查询。网页检索的数量由配置限制，默认每条断言四条查询；页面抓取只抓前若干条结果，避免不必要的网络开销。第二是锚点抽取，一次LLM调用的输入在数千token量级，输出在数百token量级。第三是读出，动作空间的大小为$|\mathcal H|+\sum_{k=2}^{|\mathcal H|}\binom{|\mathcal H|}{k}+1$，当$|\mathcal H|\le6$时最多63个动作，计算可以忽略。随机森林推断在本地完成，不产生API成本。

系统的存储成本主要来自审计文件。每条断言保存证据列表、锚点列表、信念轨迹、动作损失和LLM响应摘要。审计文件的大小通常在数百KB量级，适合长期保存和人工复核。LLM响应缓存单独存放，不进入公开审计文件。

\section{可审计性}

可审计性是FinPlanRAG的设计目标之一。系统把每个中间状态写入结构化文件，使读者可以从最终动作回溯到证据和锚点。表\ref{tab:auditfields}给出审计文件的主要字段。

\begin{table}[H]
\centering
\zihao{5}
\caption{审计文件的主要字段。}
\label{tab:auditfields}
\begin{tabular}{p{2.8cm}p{8.8cm}}
\toprule
对象 & 字段 \\
\midrule
契约 & 公司、代码、目标期间、指标、结果空间、标签定义、参考口径、截止时点 \\
证据 & 来源编号、标题、链接、发布日期、提供者、摘要、正文、是否官方来源 \\
锚点 & 类型、取值、来源编号、逐字引文、发布日期、可靠性、资格、隐含标签、理由、纠正记录 \\
信念 & 先验、仅锚点后验、融合后验、逐步更新轨迹 \\
动作 & 动作类型、点答案、裁决集、可信集、期望损失、各动作损失排序 \\
日期 & 点估计、区间、最高信念窗口、模型估计与窗口是否一致、依据 \\
\bottomrule
\end{tabular}
\end{table}

审计文件有两个用途。第一，复核。分析师可以逐条检查被采纳的锚点是否真的支持隐含标签，被拒绝的锚点是否应当被拒绝，证据的发布日期是否早于截止时点。第二，问责。当判断出现错误时，系统可以定位错误来自证据缺失、锚点抽取、资格判定、信念更新还是动作选择，而不是把错误归因于一个黑箱模型。

可审计性不等于正确性。一个锚点即使有来源和引文，也可能是对来源的误读；一个资格判定即使规则明确，也可能过严或过松。本文在案例研究中展示审计轨迹，并在第\ref{ch:conclusion}章把人工审计实验列为未来工作。当前的可审计性是设计属性，不是已经测量的用户收益。

% ==================== 第6章 实验 ====================
\chapter{实验}
\label{ch:experiment}

\section{实验设置}

\subsection{基准}

本文使用三个基准。主基准是2010年至2025年美国SEC 8-K文件中的季度营业收入指引断言，共1,005条。每条断言包含公司、目标季度、指引区间、截止时点、当前季度值、时点证据和已实现结果，标签为低于区间、区间内和高于区间。数据按时间划分为529条开发集、203条验证集和273条测试集；主结果使用227条具有XBRL核验或XBRL覆盖标签的测试断言，46条纯文本标签单独用于敏感性检验。

第二个基准是A股年报历史回测。本文选取40家非金融上市公司，覆盖消费、制造、医药、能源和科技行业，对2023年至2025年三个财年构造470条断言，其中营业收入增速232条、归母净利润增速120条、综合毛利率120条、披露日期118条。对每个财年$Y$，问题在$Y+1$年1月31日提出，此时三季报已经披露、年报尚未披露；结果来自随后披露的年报。该基准检验框架在真实中国市场上的时点判断和数值桥。

第三个基准是FinDVer组件检查。FinDVer包含700条针对美国10-K和10-Q的断言，标签为支持或反驳。本文按原数据顺序取前150条，比较固定证据预算与状态条件式证据获取。该基准评价证据获取和断言判定，不评价前瞻性结果预测，因此不与主基准的决策损失合并。

公开基准中的DocFinQA、FinQA、TAT-QA和FinanceBench的答案已经存在于文档中，Fin-Fact的标签是事实核查标签。这些基准可以检验检索、阅读和资格判定，但不包含截至时点尚未实现的结果，本文不把它们作为端到端前瞻性预测基准\citep{reddy2024docfinqa,chen2021finqa,zhu2021tatqa,islam2023financebench,rangapur2025finfact}。

\subsection{基线与基座}

本文比较三类方法。第一类是结构化基线：随机森林只使用截止时点可得的数值特征，不使用文本锚点；其训练采用公司不重叠、按年份的折划分和温度缩放。第二类是纯LLM基线：只使用锚点后验和标签分数，不使用随机森林。第三类是融合系统：锚点后验与随机森林后验按验证集选择的权重融合，并在全动作空间上读出。

{\sloppy 本文还比较六个公开检索与验证方法：FinDVer\citep{zhao2024findver}、Fin-Fact\citep{rangapur2025finfact}、Self-RAG\citep{asai2024selfrag}、CRAG\citep{yan2024crag}、FLARE\citep{jiang2023flare}和Adaptive-RAG\citep{jeong2024adaptiverag}。每个方法按原论文的判据适配到三分类结果空间，使用与本文相同的证据块、适配层和评分口径；表中省略-style后缀，六个名称均指同基座适配。\par}

三套基座为DeepSeek Flash、GLM 5.3 Flash和Qwen3-8B\citep{qwen2025qwen3}。DeepSeek和GLM通过API调用，Qwen3-8B在32GB GPU上用vLLM本地部署。所有方法使用同一基座、同一证据预算和同一标签空间。

\subsection{指标与统计检验}

本文报告四类指标。分类指标包括准确率、平衡准确率、宏F1、MCC、AUROC和AUPRC；校准指标包括期望校准误差和Brier分数；决策指标包括非对称决策损失、覆盖率、平均集合大小和弃权率；数值指标包括数值桥的平均绝对误差、平均绝对百分比误差和披露日期窗口准确率。平衡准确率是各类召回率的平均\citep{brodersen2010balanced}，宏F1是各类F1的平均\citep{sokolova2009performance}。

决策损失使用第\ref{ch:problem}章定义的参数：$\ell_{\text{over}}=12$、$\ell_{\text{under}}=8$、$\ell_{\text{wrong}}=10$、$\ell_{\text{miss}}=8$、$\ell_{\text{width}}=2$、$\ell_{\text{abstain}}=5$、$\alpha=0.1$。这些参数在测试集上不调整。置信区间由断言级配对自助法给出，重采样10,000次；准确率差异使用McNemar检验。

\subsection{训练与泄漏控制}

开发集为2010年至2018年，验证集为2019年至2021年，测试集为2022年至2025年。随机森林使用五个折种子；对每个测试年份$Y$，训练集来自其他公司折，且结果日期早于$Y$年1月1日。融合权重、温度缩放和类型库扩展只在验证集上选择，测试集不参与拟合、校准、权重选择或类型选择。所有LLM响应按请求负载的哈希缓存；API密钥只从环境变量读取，不写入结果文件。

\section{主结果}

表\ref{tab:main}按方法列出三套基座的结果。纵向是方法，横向是基座和指标，每个基座给出准确率、平衡准确率、宏F1和全动作决策损失。DeepSeek使用验证集接受的扩展类型库，GLM和Qwen保留基础类型库。

\begin{table}[htbp]
\centering
\zihao{-5}
\setlength{\tabcolsep}{2.2pt}
\caption{主结果（227条核验断言）。Acc.为准确率，Bal.为平衡准确率，Macro为宏F1，Loss为全动作决策损失。}
\label{tab:main}
\begin{tabular}{lrrrrrrrrrrrrr}
\toprule
\multirow{2}{*}{方法} & \multicolumn{4}{c}{DeepSeek Flash} & \multicolumn{4}{c}{GLM 5.3 Flash} & \multicolumn{4}{c}{Qwen3-8B} \\
\cmidrule(lr){2-5}\cmidrule(lr){6-9}\cmidrule(lr){10-13}
 & Acc. & Bal. & Macro & Loss & Acc. & Bal. & Macro & Loss & Acc. & Bal. & Macro & Loss \\
\midrule
FinDVer & .515 & .424 & .423 & 3.128 & .507 & .456 & .458 & 3.198 & .463 & .441 & .429 & 5.339 \\
Fin-Fact & .524 & .428 & .429 & 3.251 & .489 & .401 & .400 & 3.410 & .507 & .473 & .459 & 4.881 \\
Self-RAG & .546 & .460 & .468 & 3.031 & .507 & .454 & .460 & 3.181 & .485 & .446 & .431 & 5.119 \\
CRAG & .515 & .412 & .407 & 3.004 & .498 & .424 & .427 & 3.480 & .467 & .433 & .417 & 5.251 \\
FLARE & .511 & .418 & .412 & 3.278 & .489 & .407 & .408 & 3.559 & .326 & .298 & .212 & 6.661 \\
Adaptive-RAG & .489 & .402 & .396 & 3.612 & .436 & .366 & .366 & 4.070 & .352 & .325 & .290 & 6.396 \\
\midrule
LLM-only & .498 & .397 & .362 & 3.445 & .533 & .428 & .414 & 3.427 & .410 & .357 & .267 & 5.498 \\
RF & .533 & .491 & .486 & 2.987 & .533 & .491 & .486 & 2.987 & .533 & .491 & .486 & 2.987 \\
\textbf{FinPlanRAG (Ours)} & \textbf{.577} & \textbf{.526} & \textbf{.530} & \textbf{2.907} & \textbf{.551} & \textbf{.504} & \textbf{.503} & \textbf{2.943} & \textbf{.542} & \textbf{.501} & \textbf{.499} & \textbf{2.925} \\
\bottomrule
\end{tabular}
\end{table}

表中的FinPlanRAG指锚点后验与随机森林后验按验证集权重融合后的系统，融合权重只在验证集选择。DeepSeek的FinPlanRAG使用验证集扩展后的类型库，GLM和Qwen使用基础类型库；类型库扩展是训练步骤，不是另一个预测模块。FinPlanRAG在三个基座的准确率、平衡准确率、宏F1和决策损失上均为最高。DeepSeek的FinPlanRAG相对最强基线CRAG把损失从3.004降到2.907，准确率从.515提高到.577；GLM相对Self-RAG把损失从3.181降到2.943，准确率从.507提高到.551；Qwen相对Fin-Fact把损失从4.881降到2.925，准确率从.507提高到.542。配对自助法给出的损失差分别为$-0.097$ $[-0.502,0.300]$、$-0.238$ $[-0.661,0.176]$和$-1.956$ $[-2.617,-1.322]$，只有Qwen的区间排除零。

{\sloppy FinPlanRAG相对随机森林的损失差为$-0.079$ $[-0.300,0.141]$、$-0.044$ $[-0.203,0.115]$和$-0.062$ $[-0.282,0.159]$，三个区间均包含零。\par}融合在点估计上低于随机森林，但总体差异不足以主张稳定优势。把测试集按真实结果分层后，融合在区间内样本上的优势更清楚，见第\ref{sec:within}节。

表\ref{tab:calibration}给出概率质量指标。融合系统在三个基座上的AUROC为.713、.711和.707，ECE为.064、.030和.069，Brier为.555、.557和.561；对应的随机森林AUROC为.706，ECE为.033，Brier为.566。FinPlanRAG的排序能力和Brier分数优于随机森林；ECE在DeepSeek和Qwen上略高，在GLM上略低，说明融合后验的排序更准，但置信度校准仍有改进空间。

\begin{table}[htbp]
\centering
\zihao{5}
\setlength{\tabcolsep}{4pt}
\caption{概率质量与校准（227条核验断言）。AUROC和AUPRC越高越好，ECE和Brier越低越好。}
\label{tab:calibration}
\begin{tabular}{lrrrrrrrrr}
\toprule
\multirow{2}{*}{方法} & \multicolumn{3}{c}{DeepSeek Flash} & \multicolumn{3}{c}{GLM 5.3 Flash} & \multicolumn{3}{c}{Qwen3-8B} \\
\cmidrule(lr){2-4}\cmidrule(lr){5-7}\cmidrule(lr){8-10}
 & AUROC & ECE & Brier & AUROC & ECE & Brier & AUROC & ECE & Brier \\
\midrule
LLM-only & .668 & .110 & .612 & .651 & .101 & .609 & .493 & .439 & .975 \\
RF & .706 & .033 & .566 & .706 & .033 & .566 & .706 & .033 & .566 \\
\textbf{FinPlanRAG (Ours)} & \textbf{.713} & \textbf{.064} & \textbf{.555} & \textbf{.711} & \textbf{.030} & \textbf{.557} & \textbf{.707} & \textbf{.069} & \textbf{.561} \\
\bottomrule
\end{tabular}
\end{table}

\section{决策读出的作用}

表\ref{tab:readout}在同一后验上比较受限可信集门与全动作读出。受限门只比较点答案和最高密度可信集，全动作门在点答案、非单元素裁决集和弃权上选择期望损失最小的动作。融合后验在三个基座上的全动作损失相对受限门分别下降0.749、0.581和0.696，95\%配对区间均排除零。随机森林后验的下降为0.467，区间同样排除零。全动作读出没有改变后验，也没有增加训练数据，它改变的是给定后验时的动作选择。

\begin{table}[htbp]
\centering
\zihao{-5}
\setlength{\tabcolsep}{2.5pt}
\caption{受限门与全动作读出的决策损失（227条核验断言）。$\Delta$为全动作减受限门。}
\label{tab:readout}
\begin{tabular}{lrrrrrrrrr}
\toprule
\multirow{2}{*}{方法} & \multicolumn{3}{c}{DeepSeek} & \multicolumn{3}{c}{GLM} & \multicolumn{3}{c}{Qwen3} \\
\cmidrule(lr){2-4}\cmidrule(lr){5-7}\cmidrule(lr){8-10}
 & 受限门 & 全动作 & $\Delta$ & 受限门 & 全动作 & $\Delta$ & 受限门 & 全动作 & $\Delta$ \\
\midrule
\textbf{FinPlanRAG (Ours)} & \textbf{3.656} & \textbf{2.907} & \textbf{$-0.749$} & \textbf{3.524} & \textbf{2.943} & \textbf{$-0.581$} & \textbf{3.621} & \textbf{2.925} & \textbf{$-0.696$} \\
RF & 3.454 & 2.987 & $-0.467$ & 3.454 & 2.987 & $-0.467$ & 3.454 & 2.987 & $-0.467$ \\
LLM-only & 3.674 & 3.445 & $-0.229$ & 3.454 & 3.427 & $-0.026$ & 5.048 & 5.498 & $+0.449$ \\
\bottomrule
\end{tabular}
\end{table}

全动作读出的作用集中在高不确定性区域。融合后验在227条测试断言中有217条选择裁决集、10条选择点答案、0条弃权；平均集合大小为2.12。点答案不是唯一合理的输出，当两个相邻结果的期望损失接近时，返回二元裁决集可以降低误判成本。这一结果说明，前瞻性断言的核心问题不只是预测哪个标签，还包括给定后验应当输出什么。

\section{指导区间内的分层结果}
\label{sec:within}

总体损失打平时，锚点层是否在某个误差层上提供互补信号，是一个独立问题。表\ref{tab:within}报告真实结果落在指引区间内的82条测试断言。FinPlanRAG在三个基座上均低于随机森林：DeepSeek从3.146降到2.780，GLM从3.146降到2.902，Qwen从3.146降到2.732（以第一个随机森林种子为例）。五个随机森林种子中，DeepSeek和Qwen的五个区间全部排除零，GLM有四个区间排除零。

\begin{table}[htbp]
\centering
\zihao{5}
\setlength{\tabcolsep}{4pt}
\caption{指导区间内样本的分层比较（82条测试断言）。差值为FinPlanRAG减随机森林，负值表示FinPlanRAG损失更低。}
\label{tab:within}
\begin{tabular}{lrrrr}
\toprule
基座 & FinPlanRAG & 随机森林 & 损失差 [95\% CI] & 种子一致性 \\
\midrule
DeepSeek & \textbf{2.780} & 3.146 & $-0.366$ $[-0.732,-0.049]$ & 5/5 \\
GLM & \textbf{2.902} & 3.146 & $-0.244$ $[-0.537,-0.024]$ & 4/5 \\
Qwen3 & \textbf{2.732} & 3.146 & $-0.415$ $[-0.732,-0.170]$ & 5/5 \\
\bottomrule
\end{tabular}
\end{table}

这个分层结果说明锚点层在高密度、低极端误差的区域提供了互补信号。分层依据是真实结果，预测时不可知，因此它只能解释误差来源，不能作为部署时的路由规则。总体损失上的差异仍按不显著报告。

\section{A股年报历史回测}
\label{sec:ashare}

表\ref{tab:ashare}给出A股断言的分类与决策结果，DeepSeek和GLM使用470条全样本，Qwen3使用250条分层子样本。上年外推基线把上年增速或上年毛利率直接外推，季度年化基线把三季报年化后外推；完整系统使用锚点、资格判定、时间演化、数值桥和全动作读出，去桥变体关闭数值桥。

\begin{table}[htbp]
\centering
\zihao{5}
\setlength{\tabcolsep}{4pt}
\caption{A股年报历史回测。Acc.为准确率，Bal.为平衡准确率，Macro为宏F1，Loss为全动作决策损失。Qwen3在250条分层子样本上评估，其余方法在470条全样本上评估。}
\label{tab:ashare}
\begin{tabular}{llrrrr}
\toprule
方法 & 样本 & Acc. & Bal. & Macro & Loss \\
\midrule
\multicolumn{6}{l}{\textit{传统基线}} \\
上年外推 & 470 & .491 & .402 & .401 & 5.094 \\
季度年化 & 470 & .777 & .786 & .785 & 2.270 \\
\midrule
\multicolumn{6}{l}{\textit{DeepSeek}} \\
\textbf{FinPlanRAG} & 470 & \textbf{.783} & \textbf{.778} & \textbf{.776} & \textbf{2.179} \\
去数值桥 & 470 & .638 & .591 & .590 & 2.421 \\
\midrule
\multicolumn{6}{l}{\textit{GLM}} \\
\textbf{FinPlanRAG} & 470 & \textbf{.549} & \textbf{.485} & \textbf{.483} & \textbf{3.332} \\
去数值桥 & 470 & .549 & .485 & .483 & 3.332 \\
\midrule
\multicolumn{6}{l}{\textit{Qwen3}} \\
季度年化 & 250 & .772 & .789 & .789 & 2.324 \\
\textbf{FinPlanRAG} & 250 & \textbf{.580} & \textbf{.523} & \textbf{.522} & \textbf{4.208} \\
去数值桥 & 250 & .564 & .500 & .499 & 3.864 \\
\bottomrule
\end{tabular}
\end{table}

DeepSeek基座上的完整系统取得.783的准确率和2.179的决策损失，高于季度年化的.777和2.270，也明显高于去桥变体的.638和2.421。完整系统相对去桥变体的损失差为$-0.243$ $[-0.404,-0.077]$，准确率差为$+0.145$ $[0.109,0.183]$，两个区间均排除零。数值桥把营业收入和归母净利润的增速判断从标签后验推进到可计算的数值区间，是DeepSeek增益的主要来源。GLM基座上完整系统与去桥变体相同，说明该基座的抽取结果没有提供足够的数值锚点，桥接未触发；完整系统也低于季度年化。Qwen3的250条子样本上，完整系统的准确率比去桥变体高$+0.016$ $[0.004,0.032]$，但损失高$+0.344$ $[0.064,0.640]$，说明数值桥改善了部分区间判断，却没有改善整体决策校准。A股结果因此分为两类：DeepSeek上数值桥带来明确增益，GLM和Qwen上抽取质量不足以支撑桥接，标签后验仍是更稳的回退。

表\ref{tab:ashare-numeric}给出数值误差。DeepSeek完整系统的数值MAE为9.86个百分点，季度年化为7.40个百分点；完整系统的分类损失更低，但数值点估计不如简单年化。这个差异说明两种目标并不等价：数值桥改善的是区间判断和决策损失，不保证点值误差最小。披露日期窗口准确率为.788，高于季度年化的.729。Qwen3的数值MAE为39.21个百分点，日期窗口准确率为.690，说明其抽取器输出的数值字段噪声较大。

\begin{table}[H]
\centering
\zihao{5}
\setlength{\tabcolsep}{4pt}
\caption{A股回测的数值与日期指标。MAE为增速或毛利率的百分点误差，MAPE只在真实值非零时计算。}
\label{tab:ashare-numeric}
\begin{tabular}{llrrr}
\toprule
方法 & 样本 & 数值MAE & 数值MAPE & 日期窗口准确率 \\
\midrule
季度年化 & 470 & 7.40 & 459.7 & .729 \\
\midrule
\multicolumn{5}{l}{\textit{DeepSeek}} \\
\textbf{FinPlanRAG} & 470 & \textbf{9.86} & \textbf{408.7} & \textbf{.788} \\
去数值桥 & 470 & -- & -- & .788 \\
\midrule
\multicolumn{5}{l}{\textit{GLM}} \\
\textbf{FinPlanRAG} & 470 & -- & -- & \textbf{.754} \\
\midrule
\multicolumn{5}{l}{\textit{Qwen3}} \\
\textbf{FinPlanRAG} & 250 & \textbf{39.21} & \textbf{508.7} & \textbf{.690} \\
去数值桥 & 250 & -- & -- & .690 \\
\bottomrule
\end{tabular}
\end{table}

\section{公开基准的组件级检查}

表\ref{tab:findver}给出FinDVer前150条断言的结果。状态条件式获取从两条证据开始，只有在阅读器报告证据不足时才扩展，最多八条；固定预算条件直接给4条或8条；金标条件给出专家标注的证据段落。

\begin{table}[H]
\centering
\zihao{5}
\setlength{\tabcolsep}{4pt}
\caption{FinDVer组件检查（150条断言）。Acc.为断言判定准确率，Macro为宏F1，Recall为证据段落召回率，Para.为平均证据段落数。}
\label{tab:findver}
\zihao{-5}
\setlength{\tabcolsep}{2.2pt}
\begin{tabular}{lrrrrrrrrrrrr}
\toprule
\multirow{2}{*}{条件} & \multicolumn{4}{c}{DeepSeek} & \multicolumn{4}{c}{GLM} & \multicolumn{4}{c}{Qwen3} \\
\cmidrule(lr){2-5}\cmidrule(lr){6-9}\cmidrule(lr){10-13}
 & Acc. & Macro & Recall & Para. & Acc. & Macro & Recall & Para. & Acc. & Macro & Recall & Para. \\
\midrule
固定4条 & .767 & .550 & .620 & 4.00 & .760 & .485 & .620 & 4.00 & .733 & .438 & .620 & 4.00 \\
固定8条 & .847 & .570 & .755 & 8.00 & .860 & .573 & .755 & 8.00 & .700 & .426 & .755 & 8.00 \\
金标证据 & .913 & .577 & 1.000 & 2.99 & .933 & .603 & 1.000 & 2.99 & .707 & .441 & 1.000 & 2.99 \\
\textbf{FinPlanRAG 获取} & \textbf{.873} & \textbf{.559} & \textbf{.603} & \textbf{4.63} & \textbf{.913} & \textbf{.592} & \textbf{.628} & \textbf{4.95} & \textbf{.660} & \textbf{.397} & \textbf{.517} & \textbf{4.01} \\
\bottomrule
\end{tabular}
\end{table}

DeepSeek和GLM上，FinPlanRAG的状态条件式获取在约一半段落数下取得高于固定8条的判定准确率：DeepSeek为.873对.847，GLM为.913对.860。证据召回率低于固定8条，说明它没有覆盖全部金标段落，但覆盖了足以判定的段落。Qwen3上的方向相反：固定4条为.733，状态条件式为.660，低于金标条件.707。Qwen3的阅读器在短上下文下更稳定，扩展证据反而引入噪声。这个结果把状态条件式获取的适用条件限定在阅读器能够利用新增证据的基座上。

\section{误差分析}

第一类误差来自标签噪声。XBRL口径、季度划分和重述会使少量测试标签与披露文本不完全一致；46条纯文本标签不进入主表，避免把标签噪声混入比较。

第二类误差来自锚点抽取。参考值锚点剔除、日期投影和单位检查可以处理常见错误，但语义误读仍会发生。错误的硬锚点会排除正确结果，错误的软锚点会把信念拉向错误方向。GLM在A股回测中没有产生足够的数值锚点，说明抽取质量直接决定数值桥是否可用。

第三类误差来自基座能力。Qwen3-8B在长上下文断言判定上弱于DeepSeek和GLM，融合权重虽然不为零，但锚点后验的贡献有限；在FinDVer上，Qwen3从固定4条扩展到8条时准确率下降，说明新增证据没有被有效利用。三套基座的结果按基座分别报告，不合并平均值。

% ==================== 第7章 消融实验 ====================
\chapter{消融实验}
\label{ch:ablation}

本章的消融在锚点分支和最终系统上分别进行。锚点分支的损失只用于比较证据策略，不与第\ref{ch:experiment}章的融合主表直接比较；最终系统的消融使用与主表相同的后验和读出。

\section{锚点提示词与时间演化}

文本锚点的作用首先出现在纯LLM分支。旧提示词使用结构化摘要，新提示词使用时点原文、段落编号和逐字引文。DeepSeek上的准确率、平衡准确率和宏F1分别从.445、.307、.242提高到.502、.397、.364。来源落地和逐字引文让模型能够区分“指引区间”“历史实际”和“分析师预测”，减少了把累计增速当作全年增速的错误。

时间演化的作用较小但方向一致。新锚点从静态可靠度切换到按揭示日更新的Beta可靠度，损失从2.908降到2.855；旧锚点从2.899降到2.793。准确率变化不大，演化主要改善概率校准和决策损失，而不是点分类。这个结果符合演化的设计目标：它更新的是锚点类型的可靠性，不是标签本身。

\section{资格判定}

表\ref{tab:qual-ablation}比较完整资格判定与接受全部解析锚点。完整判定拒绝空泛、过期、对抗和不合格的锚点，接受率在三个基座上分别为27.6\%、42.7\%和68.3\%；接受全部条件保留所有解析结果。

\begin{table}[htbp]
\centering
\zihao{5}
\setlength{\tabcolsep}{4pt}
\caption{资格判定消融（227条核验断言，锚点分支）。差值为接受全部减完整判定，正值表示完整判定损失更低。}
\label{tab:qual-ablation}
\begin{tabular}{lrrrr}
\toprule
基座 & 完整判定 & 接受全部 & 损失差 [95\% CI] & 接受率 \\
\midrule
DeepSeek & 3.304 & 3.507 & $+0.203$ $[-0.106,0.520]$ & 27.6\% \\
GLM & 3.339 & 3.427 & $+0.088$ $[-0.176,0.361]$ & 42.7\% \\
Qwen3 & 4.599 & 4.652 & $+0.053$ $[-0.141,0.256]$ & 68.3\% \\
\bottomrule
\end{tabular}
\end{table}

三个基座的点估计都支持完整判定，DeepSeek的增益最大，Qwen3的增益最小。三个配对区间都包含零，因此资格判定的作用按方向性证据报告。Qwen3的接受率明显高于另外两个基座，说明其抽取器产生了更多边缘锚点，资格判定拒绝的比例更高，但剩余锚点的区分力有限。

\section{融合权重}

表\ref{tab:weight}给出验证集上不同融合权重的决策损失。权重$w$是锚点后验在融合后验中的占比，$w=0$对应纯随机森林，$w=1$对应纯锚点后验。DeepSeek在$w=0.2$取得最低验证损失，GLM和Qwen3在$w=0.1$取得最低验证损失；测试集使用验证集选择的权重，不参与调参。

\begin{table}[htbp]
\centering
\zihao{5}
\setlength{\tabcolsep}{4pt}
\caption{融合权重的验证集损失（185条验证断言）。粗体为验证集选择。}
\label{tab:weight}
\begin{tabular}{lrrrrrrr}
\toprule
基座 & 0.0 & 0.05 & 0.1 & 0.2 & 0.3 & 0.5 & 1.0 \\
\midrule
DeepSeek & 3.341 & 3.308 & 3.265 & \textbf{3.232} & 3.341 & 3.503 & 3.859 \\
GLM & 3.341 & 3.211 & \textbf{3.200} & 3.286 & 3.362 & 3.416 & 3.708 \\
Qwen3 & 3.341 & 3.362 & \textbf{3.330} & 3.686 & 3.859 & 4.162 & 6.519 \\
\bottomrule
\end{tabular}
\end{table}

三个基座的最优权重都偏向随机森林，DeepSeek为0.2，GLM和Qwen3为0.1。锚点后验单独使用时损失明显更高，Qwen3在$w=1$时达到6.519。融合的价值来自补充信号，而不是替代结构化后验。

\section{类型库的验证集扩展}

类型库扩展是算法的一个训练步骤，不是独立的预测模块。算法\ref{alg:typeexpand}先在验证集上挖掘高损失断言，由LLM提出候选类型，再在选择集上重抽锚点并计算损失；只有验证损失下降超过$\delta$的候选才进入类型库。接受后的类型库冻结，测试集只评估一次。


DeepSeek上的候选类型库包含“指引中点与当前季度收入之比”和“财年周数效应”两个类型。验证选择集上，基础类型库的损失为3.943，候选类型库为3.792，下降0.151，超过阈值，候选被接受。提示词单独修改和类型加提示词组合两个候选没有通过验证门控。冻结类型库后，测试集上的准确率从.555提高到.577，宏F1从.513提高到.530，损失从2.934降到2.907；配对损失差为$-0.026$ $[-0.159,0.097]$，区间包含零。

表\ref{tab:typeevolve}汇总这一过程。类型库扩展的测试增益是小的，但方向与验证集一致，且没有使用测试标签。GLM和Qwen3没有候选通过验证门控，保留基础类型库，因此主表中只有DeepSeek使用扩展类型库。

\begin{table}[htbp]
\centering
\zihao{5}
\setlength{\tabcolsep}{4pt}
\caption{DeepSeek类型库扩展（227条核验断言）。Loss为全动作决策损失。}
\label{tab:typeevolve}
\begin{tabular}{lrrrr}
\toprule
系统 & Acc. & Bal. & Macro & Loss \\
\midrule
基础类型库 & .555 & .511 & .513 & 2.934 \\
扩展类型库 & .577 & .526 & .530 & 2.907 \\
差值 & $+.022$ & $+.015$ & $+.017$ & $-0.026$ $[-0.159,0.097]$ \\
\bottomrule
\end{tabular}
\end{table}

\section{去掉锚点结构的对照}
\label{sec:no-anchor}

前面的消融检验的是组件内部，例如资格判定是否比接受全部更好。一个更根本的问题是：如果去掉三层锚点结构，只保留时点证据、LLM标签分数、温度缩放和全动作读出，损失会差多少。这个对照直接检验锚点结构本身的必要性。

本文从同一批锚点抽取响应中只保留标签分数，去掉全部锚点对象、资格判定和Beta可靠性，在验证集上拟合温度参数，再在测试集上使用同一个全动作读出。这个对照可以看作FinAbstain式选择性预测在本任务上的近似：时点证据、概率输出、温度缩放和损失感知弃权都在，但没有多类别资格判定和类型可靠性演化。表\ref{tab:noanchor}给出结果。去掉锚点结构后，DeepSeek、GLM和Qwen的损失分别上升0.379、0.317和0.899，准确率分别下降5.3、8.8和13.7个百分点。时点证据和选择性输出仍然存在，性能却明显低于完整系统，说明增益不能只归因于“时点检索加弃权”这一一般性组合。

\begin{table}[htbp]
\centering
\zihao{5}
\setlength{\tabcolsep}{4pt}
\caption{去掉三层锚点结构的架构级对照（227条核验断言）。温度在验证集上选择，测试集只评估一次。}
\label{tab:noanchor}
\begin{tabular}{lrrrrrrrrr}
\toprule
\multirow{2}{*}{方法} & \multicolumn{3}{c}{DeepSeek} & \multicolumn{3}{c}{GLM} & \multicolumn{3}{c}{Qwen3} \\
\cmidrule(lr){2-4}\cmidrule(lr){5-7}\cmidrule(lr){8-10}
 & Acc. & Macro & Loss & Acc. & Macro & Loss & Acc. & Macro & Loss \\
\midrule
无锚点结构 & .524 & .398 & 3.286 & .463 & .347 & 3.260 & .405 & .257 & 3.824 \\
\textbf{FinPlanRAG} & \textbf{.577} & \textbf{.530} & \textbf{2.907} & \textbf{.551} & \textbf{.504} & \textbf{2.943} & \textbf{.542} & \textbf{.501} & \textbf{2.925} \\
\bottomrule
\end{tabular}
\end{table}

这个对照也解释了为什么本文把贡献放在证据资格判定和类型可靠性演化上，而不是放在“选择性预测”上。选择性预测本身是一个一般性工具，FinAbstain等工作已经把它用于金融预测；本文的增量在于锚点结构把证据拆成可判定、可追溯、可按类型更新的对象。去掉这层结构，系统仍然可以选择弃权或裁决集，但失去了区分证据失效模式的能力。

\section{数值桥}

数值桥的作用在A股回测上最清楚。DeepSeek完整系统的损失为2.179，去桥变体为2.421，差值为$-0.242$；准确率从.638提高到.783，宏F1从.590提高到.776。桥接把“增速可能高于区间”的标签判断转化为“按当前季度年化，全年增速约为多少”的数值计算，减少了纯标签后验的过度平滑。

桥接的收益依赖抽取质量。GLM完整系统与去桥变体完全相同，说明该基座没有产生足够的数值锚点，桥接未触发；Qwen3的结果见第\ref{ch:experiment}章。数值桥不是无条件增益，它要求抽取器输出可核验的数值字段、单位和期间。对抽取质量较弱的基座，标签后验仍是更稳的回退。

\section{消融结论}

消融结果可以归纳为四点。第一，去掉三层锚点结构后，三个基座的损失分别上升0.379、0.317和0.899，说明锚点结构本身提供了不可替代的增量，而不是只有选择性输出在起作用。第二，输出形式和资格判定是稳定的组件：全动作读出在三个基座上显著降低损失，资格判定方向一致但区间包含零。第三，类型库扩展和数值桥是条件增益：它们在DeepSeek上有效，在抽取质量较弱的基座上不触发或收益有限。第四，融合权重偏向随机森林，锚点层提供的是补充信号，而不是替代结构化后验。这个结论与主表的定位一致：框架的优势来自证据状态、决策接口和可审计性，而不是单一组件的普遍提升。

% ==================== 第8章 讨论与结论 ====================
\chapter{讨论与结论}
\label{ch:conclusion}

\section{讨论}

本文的实验支持一个分层结论，也决定了论文的主张边界。决策接口是稳健的：融合后验的全动作读出在三个基座上相对受限门分别降低0.749、0.581和0.696的决策损失，随机森林后验降低0.467，融合与随机森林共六个配对区间全部排除零。证据状态是有竞争力的：融合系统在三个基座的准确率、平衡准确率、宏F1和决策损失上均为最高，相对最强基线的点估计更低，但只有Qwen的区间排除零。误差分层提供了更清楚的证据：在真实结果落在指引区间内的82条断言上，三个基座的融合损失都显著低于随机森林，DeepSeek和Qwen的五个种子全部排除零，GLM有四个种子排除零。

这个分层对金融应用有意义。一个总体损失与强结构化基线相当的系统，仍然可以在特定误差层上提供互补信号，并且可以把证据、资格、可靠性和动作完整保留下来。金融决策关心的不只是平均准确率，还包括错误发生在哪个区域、错误成本是否对称、输出是否可复核。本文把这些维度放进同一个契约和同一个损失，因而能够分别报告后验质量、动作空间和证据状态的作用。

因此本文不把全面性能优越作为唯一主张。锚点层的作用应定位为可审计的语义状态和证据资格判定协议，而不是普遍的点预测优势。文本锚点把纯LLM分支的准确率从.445提高到.502，融合后验在区间内样本上显著优于随机森林，但在总体损失上与随机森林的差异不显著。这个结果不削弱框架的价值：监管、审计和研究场景需要的正是一个能够说明“为什么这样判断”的状态，而不是一个只返回标签的黑箱。

外部基准给出了适用边界。FinDVer上，状态条件式获取在DeepSeek和GLM上以约一半的段落数取得高于固定8条的判定准确率，在Qwen3上则低于固定4条。A股回测上，数值桥在DeepSeek上明显改善损失，在GLM上没有触发。两个结果指向同一个条件：证据策略和数值桥的收益取决于基座的抽取能力。框架提供接口，基座决定接口能否被充分利用。

\section{局限}

第一，主基准只覆盖美国公司的季度营业收入指引。A股回测补充了中国市场的年报断言，但两个市场的制度、披露习惯和行业结构不同，跨市场结论仍需更多数据。FinDVer只用于组件检查，不构成端到端前瞻性预测的外部验证。

第二，标签来自XBRL核验或披露文件回退，不是人工金标准。季度口径、重述和单位差异会带来噪声；46条纯文本标签被排除在主结果之外，缩小了测试集。A股回测的数值桥依赖抽取器输出数值字段，GLM上的桥接未触发，Qwen的数值误差较大，说明抽取质量仍是实际部署的主要瓶颈。

第三，类型库扩展只在DeepSeek上完成一轮，GLM和Qwen3没有候选通过验证门控。单轮扩展的测试增益较小，配对区间包含零；更多轮次、跨基座类型迁移和人工审计留待后续工作。

\section{结论}

本文把前瞻性金融断言表示为截止时点上的契约，并在有限结果空间上维护可演化的锚点状态。框架由三层锚点、五类资格判定、揭示日Beta可靠性、验证集门控的类型库扩展、确定性数值桥和全动作损失感知读出组成。实验在227条核验断言、470条A股年报断言和150条FinDVer断言上进行，覆盖三个基座。

结果支持三点结论。第一，全动作读出是稳健的决策改进，在三个基座上显著降低决策损失。第二，融合系统在三个基座的分类和决策指标上均为最高，并在指导区间内样本上显著优于随机森林。第三，类型库扩展和数值桥是条件增益，其效果取决于基座的抽取质量。框架的主要贡献不是某个单点的准确率提升，而是把证据、判断和动作组织成一个可检查、可演化、可复用的决策接口。

\clearpage
\phantomsection
\addcontentsline{toc}{chapter}{参考文献}
{\zihao{5}\setlength{\baselineskip}{16pt}\begin{thebibliography}{28}
\setlength{\bibsep}{0pt}
\setlength{\itemsep}{0pt}
\setlength{\parsep}{0pt}
\setlength{\bibhang}{2em}

\bibitem[Angelopoulos and Bates(2021)]{angelopoulos2021gentle}
Angelopoulos, A. N., and Bates, S., 2021, ``A Gentle Introduction to Conformal Prediction and Distribution-Free Uncertainty Quantification'', arXiv:2107.07511.

\bibitem[Asai et al.(2024)]{asai2024selfrag}
Asai, A., Wu, Z., Wang, Y., Sil, A., and Hajishirzi, H., 2024, ``Self-RAG: Learning to Retrieve, Generate, and Critique through Self-Reflection'', in \emph{International Conference on Learning Representations}.

\bibitem[Brodersen et al.(2010)]{brodersen2010balanced}
Brodersen, K. H., Ong, C. S., Stephan, K. E., and Buhmann, J. M., 2010, ``The Balanced Accuracy and Its Posterior Distribution'', in \emph{Proceedings of the 20th International Conference on Pattern Recognition}, pp. 3121--3124.

\bibitem[Chen et al.(2021)]{chen2021finqa}
Chen, Z., Chen, W., Smiley, C., Shah, S., Borova, I., Langdon, D., Moussa, R., Beane, M., Huang, T.-H., Routledge, B., and Wang, W. Y., 2021, ``FinQA: A Dataset of Numerical Reasoning over Financial Data'', in \emph{Proceedings of the 2021 Conference on Empirical Methods in Natural Language Processing}, pp. 3697--3711.

\bibitem[Chow(1970)]{chow1970optimum}
Chow, C. K., 1970, ``On Optimum Recognition Error and Reject Tradeoff'', \emph{IEEE Transactions on Information Theory}, Vol. 16, No. 1, pp. 41--46.

\bibitem[Elkan(2001)]{elkan2001cost}
Elkan, C., 2001, ``The Foundations of Cost-Sensitive Learning'', in \emph{Proceedings of the 17th International Joint Conference on Artificial Intelligence}, pp. 973--978.

\bibitem[Geifman and El-Yaniv(2017)]{geifman2017selective}
Geifman, Y., and El-Yaniv, R., 2017, ``Selective Classification for Deep Neural Networks'', in \emph{Advances in Neural Information Processing Systems}, Vol. 30.

\bibitem[Gibbs and Cand\`es(2021)]{gibbs2021adaptive}
Gibbs, I., and Cand\`es, E. J., 2021, ``Adaptive Conformal Inference Under Distribution Shift'', in \emph{Advances in Neural Information Processing Systems}, Vol. 34, pp. 1660--1672.

\bibitem[Guo et al.(2017)]{guo2017calibration}
Guo, C., Pleiss, G., Sun, Y., and Weinberger, K. Q., 2017, ``On Calibration of Modern Neural Networks'', in \emph{Proceedings of the 34th International Conference on Machine Learning}, Vol. 70, pp. 1321--1330.

\bibitem[Hirst et al.(2008)]{hirst2008management}
Hirst, D. E., Koonce, L., and Venkataraman, S., 2008, ``Management Earnings Forecasts: A Review and Framework'', \emph{Accounting Horizons}, Vol. 22, No. 3, pp. 315--338.

\bibitem[Islam et al.(2023)]{islam2023financebench}
Islam, P., Kannappan, A., Kiela, D., Qian, R., Scherrer, N., and Vidgen, B., 2023, ``FinanceBench: A New Benchmark for Financial Question Answering'', arXiv:2311.11944.

\bibitem[Jeong et al.(2024)]{jeong2024adaptiverag}
Jeong, S., Baek, J., Cho, S., Hwang, S. J., and Park, J. C., 2024, ``Adaptive-RAG: Learning to Adapt Retrieval-Augmented Large Language Models through Question Complexity'', in \emph{Proceedings of the 2024 Conference of the North American Chapter of the Association for Computational Linguistics}, pp. 7036--7050.

\bibitem[Jiang et al.(2023)]{jiang2023flare}
Jiang, Z., Xu, F., Gao, L., Sun, Z., Liu, Q., Dwivedi-Yu, J., Yang, Y., Callan, J., and Neubig, G., 2023, ``Active Retrieval Augmented Generation'', in \emph{Proceedings of the 2023 Conference on Empirical Methods in Natural Language Processing}, pp. 7969--7992.

\bibitem[Lewis et al.(2020)]{lewis2020rag}
Lewis, P., Perez, E., Piktus, A., Petroni, F., Karpukhin, V., Goyal, N., Kuttler, H., Lewis, M., Yih, W.-T., Rockt\"aschel, T., Riedel, S., and Kiela, D., 2020, ``Retrieval-Augmented Generation for Knowledge-Intensive NLP Tasks'', in \emph{Advances in Neural Information Processing Systems}, Vol. 33, pp. 9459--9474.

\bibitem[L\'opez de Prado(2018)]{lopezdeprado2018advances}
L\'opez de Prado, M., 2018, \emph{Advances in Financial Machine Learning}, Wiley.

\bibitem[Manski(2003)]{manski2003partial}
Manski, C. F., 2003, \emph{Partial Identification of Probability Distributions}, Springer.

\bibitem[Qwen Team(2025)]{qwen2025qwen3}
Qwen Team, 2025, ``Qwen3 Technical Report'', arXiv:2505.09388.

\bibitem[Rangapur et al.(2025)]{rangapur2025finfact}
Rangapur, A., Wang, H., Jian, L., and Shu, K., 2025, ``Fin-Fact: A Benchmark Dataset for Multimodal Financial Fact-Checking and Explanation Generation'', in \emph{Companion Proceedings of the ACM Web Conference 2025}, pp. 785--788.

\bibitem[Reddy et al.(2024)]{reddy2024docfinqa}
Reddy, V., Koncel-Kedziorski, R., Lai, V. D., Krumdick, M., Lovering, C., and Tanner, C., 2024, ``DocFinQA: A Long-Context Financial Reasoning Dataset'', in \emph{Proceedings of the 62nd Annual Meeting of the Association for Computational Linguistics (Volume 2: Short Papers)}, pp. 445--458.

\bibitem[Romano et al.(2020)]{romano2020classification}
Romano, Y., Sesia, M., and Cand\`es, E. J., 2020, ``Classification with Valid and Adaptive Coverage'', in \emph{Advances in Neural Information Processing Systems}, Vol. 33, pp. 3581--3591.

\bibitem[Shinn et al.(2023)]{shinn2023reflexion}
Shinn, N., Cassano, F., Berman, A., Gopinath, A., Narasimhan, K., and Yao, S., 2023, ``Reflexion: Language Agents with Verbal Reinforcement Learning'', in \emph{Advances in Neural Information Processing Systems}, Vol. 36.

\bibitem[Sokolova and Lapalme(2009)]{sokolova2009performance}
Sokolova, M., and Lapalme, G., 2009, ``A Systematic Analysis of Performance Measures for Classification Tasks'', \emph{Information Processing and Management}, Vol. 45, No. 4, pp. 427--437.

\bibitem[Vovk et al.(2005)]{vovk2005algorithmic}
Vovk, V., Gammerman, A., and Shafer, G., 2005, \emph{Algorithmic Learning in a Random World}, Springer.

\bibitem[Xie et al.(2024)]{xie2024finben}
Xie, Q., Han, W., Chen, Z., Xiang, R., Zhang, X., He, Y., Xiao, M., Li, D., Dai, Y., Feng, D., Xu, Y., Kang, H., Kuang, Z., Yuan, C., Yang, K., Luo, Z., Zhang, T., Liu, Z., Xiong, G., Deng, Z., Jiang, Y., Yao, Z., Li, H., Yu, Y., Hu, G., Huang, J., Liu, X.-Y., Lopez-Lira, A., Wang, B., Lai, Y., Wang, H., Peng, M., Ananiadou, S., and Huang, J., 2024, ``FinBen: A Holistic Financial Benchmark for Large Language Models'', in \emph{Advances in Neural Information Processing Systems}, Vol. 37.

\bibitem[Yan et al.(2024)]{yan2024crag}
Yan, S.-Q., Gu, J.-C., Zhu, Y., and Ling, Z.-H., 2024, ``Corrective Retrieval Augmented Generation'', arXiv:2401.15884.

\bibitem[Yao et al.(2023)]{yao2023react}
Yao, S., Zhao, J., Yu, D., Du, N., Shafran, I., Narasimhan, K., and Cao, Y., 2023, ``ReAct: Synergizing Reasoning and Acting in Language Models'', in \emph{International Conference on Learning Representations}.

\bibitem[Zhao et al.(2024)]{zhao2024findver}
Zhao, Y., Long, Y., Jiang, T., Wang, C., Chen, W., Liu, H., Tang, X., Zhang, Y., Zhao, C., and Cohan, A., 2024, ``FinDVer: Explainable Claim Verification over Long and Hybrid-Content Financial Documents'', in \emph{Proceedings of the 2024 Conference on Empirical Methods in Natural Language Processing}, pp. 14739--14752.

\bibitem[Zhu et al.(2021)]{zhu2021tatqa}
Zhu, F., Lei, W., Huang, Y., Wang, C., Zhang, S., Lv, J., Feng, F., and Chua, T.-S., 2021, ``TAT-QA: A Question Answering Benchmark on a Hybrid of Tabular and Textual Content in Finance'', in \emph{Proceedings of the 59th Annual Meeting of the Association for Computational Linguistics}, pp. 3277--3287.


\bibitem[Guan et al.(2026)]{guan2026finverify}
Guan, M., Chen, Q., Meng, G., and Liu, X.-Y., 2026, ``FinVerify: A Verification Framework for Financial Agents'', arXiv:2611.20423.

\bibitem[He et al.(2026)]{he2026chrono}
He, S., Lv, Y., Manela, A., and Wu, J., 2026, ``Chronologically Consistent Large Language Models'', \emph{Journal of Financial Economics}, Vol. 185, 104374.

\bibitem[Joren et al.(2025)]{joren2025sufficient}
Joren, H., Zhang, J., Ferng, C.-S., Juan, D.-C., Taly, A., and Rashtchian, C., 2025, ``Sufficient Context: A New Lens on Retrieval Augmented Generation Systems'', in \emph{International Conference on Learning Representations}.

\bibitem[Karger et al.(2025)]{karger2025forecastbench}
Karger, E., Bastani, H., Yueh-Han, C., Jacobs, Z., Halawi, D., Zhang, F., and Tetlock, P. E., 2025, ``ForecastBench: A Dynamic Benchmark of AI Forecasting Capabilities'', in \emph{International Conference on Learning Representations}.

\bibitem[Mao et al.(2025)]{mao2025selectllm}
Mao, Y., Durand, A., Mehrasa, N., He, Y., and Ester, M., 2025, ``Calibrating LLMs for Selective Prediction: Balancing Coverage and Risk'', in \emph{Advances in Neural Information Processing Systems}, Vol. 38.

\bibitem[Tao et al.(2026)]{tao2026tempfinrag}
Tao, L., Li, Z., Ji, Y., Liao, C.-T., Xiao, X., and Cao, X., 2026, ``TempFinRAG: Multimodal Temporal Retrieval-Augmented Generation for Point-in-Time Financial Question Answering'', \emph{Symmetry}, Vol. 18, No. 9, 1498.

\bibitem[Yin et al.(2026)]{yin2026epsbench}
Yin, Y., Xu, C., Ye, L., Peng, J., and Zhang, C., 2026, ``EPSBench: Benchmarking LLM Agents for Earnings Forecasting'', arXiv:2610.38293.

\bibitem[Torres et al.(2026)]{torres2026finabstain}
Torres, D., Cheng, W., and Huang, H., 2026, ``FinAbstain: Uncertainty-Calibrated Multimodal RAG for Selective Financial Forecasting'', arXiv:2607.24875.

\bibitem[Wang et al.(2026)]{wang2026ane}
Wang, Z., Zeng, Y., Gong, Z., Guo, Y., Zhu, F., Zhang, H., Zhang, W., and Zuo, W., 2026, ``AnE: Pushing the Reasoning Frontier of Multimodal LLMs via Anchor Evolution'', in \emph{European Conference on Computer Vision (ECCV)}, LNCS Vol. 17016, pp. 410--428.

\end{thebibliography}
}

\appendix

\chapter{锚点抽取提示词}
\label{app:prompt}

本节及附录 F 中的仓库路径均相对于 GitHub 仓库\url{\finplanrepo}。本节给出锚点抽取提示词的要点，完整提示词位于\repolink{experiments/anchor_cli/prompts/anchor_general.txt}{\texttt{experiments/anchor\_cli/prompts/anchor\_general.txt}}。

提示词的第一部分给出断言契约，包括公司、代码、目标期间、指标、截止时点、参考口径、结果标签及其定义。标签定义必须互斥且穷尽，例如“低于140\%”“140\%至170\%”“高于170\%”。提示词明确说明，指引区间定义三个标签，但区间本身不是“结果落在区间内”的证据。

第二部分给出时点证据。每条证据带来源编号$W\#$、标题、发布日期、链接和摘录。对于申报文件，摘录带段落编号$P\#$；对于结构化证据，编号为$E\#$。提示词要求锚点只引用这些编号，不引用不存在的来源。

第三部分给出任务和规则。规则包括：每条锚点必须引用一个来源编号；引文必须逐字复制；发布日期不得晚于截止时点；可靠性是隐含标签正确的校准概率；硬锚点必须是显式约束；监管截止日是上界，不是晚披露的证据；引文中提到的日期不得晚于来源发布日期；证据不足时在缺失字段中说明；标签分数必须使用契约标签且和为1。

第四部分给出输出格式。输出是一个JSON对象，包含锚点列表、标签分数、日期估计（仅日期断言）、置信度和缺失说明。锚点字段包括类型、关系、取值、作用域、来源编号、引文、日期、可靠性、硬标记、隐含标签和理由。日期估计字段包括点估计、区间下界、区间上界和依据。

提示词的设计目标是可检查性，而不是最大化抽取数量。系统宁可返回较少的锚点，也不接受没有来源或没有引文的锚点。解析失败的响应被记录为错误，不进入信念更新；系统不使用金标结果作为回退。

\chapter{数据构造与标签核验}
\label{app:data}

收入指引断言来自2010年至2025年美国SEC 8-K文件。每条断言包含以下字段：公司代码、公司名称、行业、目标年份、目标季度、指引下限、指引中点、指引上限、指引宽度、当前季度值、当前季度可靠性、截止时点、证据项、已实现结果和结果标签。

指引区间从8-K的业绩发布文本中抽取，单位为百万美元。当前季度值来自同一份8-K或最近的10-Q。已实现结果来自SEC XBRL；XBRL没有独立季度数据时，使用披露文件回退。结果标签按实际值相对指引区间的位置确定：低于下限为“低于”，落在区间内为“区间内”，高于上限为“高于”。

标签核验分为两级。第一级是XBRL核验：实际值来自XBRL的季度事实，且单位、期间和公司代码一致。第二级是XBRL覆盖：实际值来自XBRL覆盖的披露文件，但季度划分需要回退。227条测试断言满足其中一级；46条只有纯文本标签，被排除在主结果之外，用于敏感性检验。

A股回测覆盖40家非金融上市公司和2023年至2025年三个财年，共470条断言。财务数值来自akshare的公开财务报表，公告日期和链接来自巨潮资讯网。对每个财年$Y$，截止时点为$Y+1$年1月31日；营业收入增速、归母净利润增速和综合毛利率的阈值由上年实际值加减固定带宽得到，披露日期按3月31日和4月15日划分为三个窗口。只有上年值、三季报值、当年实际值和公告记录都存在的样本才进入基准。

数据的已知限制包括：标签来自自动核验而非人工金标准；类别不平衡；主基准只覆盖美国公司的季度营业收入指引；指引区间由公司选择，区间宽度与信息环境相关；部分证据项的日期来自文件提交日，可能与实际披露时点相差一天。这些限制在第\ref{ch:conclusion}章讨论。

\chapter{审计报告示例}
\label{app:audit}

本节给出中际旭创案例的审计报告结构。审计报告的生成方式与字段说明见\repolink{docs/anchor_cli.md}{\texttt{docs/anchor\_cli.md}}；运行后报告保存在本地结果目录的\texttt{audit.md}和\texttt{audit.json}中。

审计报告的第一部分是运行元数据：运行编号、生成时间、截止时点、问题、公司、代码、模型和配置。第二部分是结论摘要：每条断言的动作、点答案、裁决集、期望损失和披露日期估计。第三部分是断言契约：指标、类型、参考口径、结果标签及其定义。

第四部分是时点证据表：编号、日期、来源、标题、链接、检索词。官方公告记录排在前面，网页结果排在后面。第五部分是锚点与资格判定表：类型、取值或引文、来源、日期、引文核验、可靠性、资格、隐含标签和理由。被拒绝的锚点保留在表中，理由逐条列出。

第六部分是信念更新表：每个标签的先验、仅锚点后验和融合后验。第七部分是全动作损失感知读出：动作、点答案、裁决集、可信集、期望损失和各动作损失排序。第八部分是披露日期估计：点估计、区间、最高信念窗口、模型估计与窗口是否一致、依据。

审计报告的最后是审计说明：每条被采纳的锚点都有来源、引文、日期、资格和可靠性；每条被拒绝的锚点都有理由；晚于截止时点的证据被列入剔除列表；引文核验采用规范化后的逐字匹配；未通过核验的锚点被降权且不能成为硬锚点。报告明确说明，它是时点判断，不是对已实现结果的复核。

\chapter{关键比较的配对区间}
\label{app:paired}

表\ref{tab:paired}汇总正文关键比较的配对损失差与95\%置信区间。差值定义为前者减后者，负值表示前者损失更低。区间使用10,000次断言级配对自助法。

\begin{table}[htbp]
\centering
\zihao{5}
\setlength{\tabcolsep}{4pt}
\caption{关键比较的配对损失差与95\%置信区间（227条核验断言）。}
\label{tab:paired}
\begin{tabular}{llrr}
\toprule
比较 & 基座 & $\Delta$损失 & 95\%置信区间 \\
\midrule
全动作读出 $-$ 受限门 & DeepSeek & $-0.749$ & $[-1.075,\,-0.405]$ \\
 & GLM & $-0.581$ & $[-0.899,\,-0.264]$ \\
 & Qwen3-8B & $-0.696$ & $[-0.987,\,-0.405]$ \\
\midrule
融合 $-$ 随机森林 & DeepSeek & $-0.079$ & $[-0.300,\,0.141]$ \\
 & GLM & $-0.044$ & $[-0.203,\,0.115]$ \\
 & Qwen3-8B & $-0.062$ & $[-0.282,\,0.159]$ \\
\midrule
融合 $-$ 随机森林 & DeepSeek，指导区间内 & $-0.366$ & $[-0.732,\,-0.049]$ \\
 & GLM，指导区间内 & $-0.244$ & $[-0.537,\,-0.024]$ \\
 & Qwen3-8B，指导区间内 & $-0.415$ & $[-0.732,\,-0.170]$ \\
\bottomrule
\end{tabular}
\end{table}

\chapter{案例研究：中际旭创2023与2024年年报}
\label{app:case}

本案例使用已经发生的历史结果。对2023财年，截止时点为2024年1月31日，此时2023年三季报已经披露、2023年年报尚未披露；证据包括2021年和2022年年报、2023年三季报和2022年三季报。2023年年报于2024年4月22日披露，实际值用于验证预测。对2024财年，截止时点为2025年1月31日，此时2024年三季报已经披露、2024年年报尚未披露；2024年三季报和2023年年报进入证据，2024年年报于2025年4月21日披露。两个截止时点使用同一套契约、阈值规则、锚点抽取提示词和全动作读出。

表\ref{tab:case2023}给出2023年的运行结果。四条断言中，营业收入、综合毛利率和披露日期的点答案正确，归母净利润的点答案错误；三个点答案和三个裁决集覆盖了真实结果，归母净利润的真实值高于54.57\%，落在裁决集之外。归母净利润的数值桥估计为41.18\%，实际为77.58\%，误差主要来自2023年第四季度光模块需求的超预期增长，历史区间和前三季度年化都没有覆盖这一变化。这个错误保留在案例中，用来说明数值桥在需求突变时的失效方式。四条断言的后验按标签顺序分别为：营业收入 .607/.327/.066，归母净利润 .603/.331/.066，综合毛利率 .045/.420/.536，披露日期 .188/.223/.588。

\begin{table}[htbp]
\centering
\zihao{-5}
\setlength{\tabcolsep}{2pt}
\caption{中际旭创2023年年报案例（截止2024-01-31）。}
\label{tab:case2023}
\begin{tabular}{>{\raggedright\arraybackslash}p{3.0cm}>{\raggedright\arraybackslash}p{1.1cm}>{\raggedright\arraybackslash}p{2.6cm}>{\raggedright\arraybackslash}p{4.6cm}>{\raggedright\arraybackslash}p{2.0cm}}
\toprule
断言 & 动作 & 点答案 & 裁决集 & 实际 \\
\midrule
营业收入同比增速 & 裁决 & 低于15.29\% & 低于15.29\%, 15.29\%至35.29\% & 11.16\% \\
归母净利润同比增速 & 裁决 & 低于24.57\% & 低于24.57\%, 24.57\%至54.57\% & 77.58\% \\
综合毛利率 & 裁决 & 高于30.31pp & 28.31pp至30.31pp, 高于30.31pp & 32.99\% \\
披露日期 & 裁决 & 晚于2024-04-15 & 2024-04-01至04-15, 晚于04-15 & 2024-04-22 \\
\bottomrule
\end{tabular}
\end{table}

表\ref{tab:case2024}给出2024年的运行结果。2023年年报披露后，2023年实际值成为2024年断言的上年实际值，历史窗口从2021至2023年移动到2022至2024年，2024年三季报进入最新证据。四条断言的集合全部覆盖真实结果，点答案也全部正确。营业收入和归母净利润的数值桥估计分别为115.38\%和130.22\%，实际为122.64\%和137.93\%，误差比2023年明显缩小；综合毛利率的点答案为31.99pp至33.99pp，实际为33.80\%。

\begin{table}[htbp]
\centering
\zihao{-5}
\setlength{\tabcolsep}{2pt}
\caption{中际旭创2024年年报案例（截止2025-01-31）。}
\label{tab:case2024}
\begin{tabular}{>{\raggedright\arraybackslash}p{3.0cm}>{\raggedright\arraybackslash}p{1.1cm}>{\raggedright\arraybackslash}p{2.6cm}>{\raggedright\arraybackslash}p{4.6cm}>{\raggedright\arraybackslash}p{2.0cm}}
\toprule
断言 & 动作 & 点答案 & 裁决集 & 实际 \\
\midrule
营业收入同比增速 & 裁决 & 高于21.16\% & 低于1.162\%, 高于21.16\% & 122.64\% \\
归母净利润同比增速 & 裁决 & 高于92.58\% & 低于62.58\%, 高于92.58\% & 137.93\% \\
综合毛利率 & 点答 & 31.99pp至33.99pp & 31.99pp至33.99pp & 33.80\% \\
披露日期 & 裁决 & 晚于2025-04-15 & 2025-04-01至04-15, 晚于04-15 & 2025-04-21 \\
\bottomrule
\end{tabular}
\end{table}

{\sloppy 四条断言的后验按标签顺序分别为：营业收入 .188/.114/.698，归母净利润 .189/.117/.694，综合毛利率 .054/.818/.129，披露日期 .185/.220/.595。\par}

两个年度的对照展示了框架的演化过程。2023年结果揭示后，已实现值进入下一轮的上年实际值和历史窗口，锚点类型可靠性按揭示日更新，2024年断言的后验和动作随之变化。演化不是重写测试标签，而是把已经公开的结果放进下一时点的证据状态。2023年归母净利润的错误没有被掩盖，2024年同一指标的集合覆盖和点答案都得到修正；这比只展示成功案例更能说明框架在真实数据上的行为和边界。

案例的证据来自巨潮资讯网的年度报告和季度报告公告，财务数值来自公开财务报表。审计文件记录每条证据的来源、日期、引文、锚点资格和动作损失。案例的截止时点、阈值和标签定义在运行前固定，实际结果只在预测完成后用于验证。

\chapter{本地命令行与Claude Code接口}
\label{app:cli}

{\sloppy 本文实现了一个本地命令行工具和一个MCP服务。命令行工具为\repolink{scripts/anchor_cli.py}{\texttt{scripts/anchor\_cli.py}}，MCP服务为\repolink{scripts/anchor_mcp.py}{\texttt{scripts/anchor\_mcp.py}}，仓库根目录的\repolink{.mcp.json}{\texttt{.mcp.json}}注册了\texttt{finplan-anchor}服务。终端运行和Claude Code运行调用同一个引擎，生成同样的审计文件。\par}

命令行提供五个命令。\texttt{case}运行\texttt{configs/anchor\_cases}下的预置案例；\texttt{ask}把自然语言问题解析成断言契约并运行；\texttt{search}只检索证据，不做推断；\texttt{audit}打印某次运行的审计报告；\texttt{serve}启动MCP服务。典型用法为：

\begin{lstlisting}
git clone https://github.com/Topps-2025/FinPlanRAG.git
cd FinPlanRAG
python scripts/anchor_cli.py ask "中际旭创2023年年报增速区间" \
    --company 中际旭创 --ticker 300308.SZ \
    --period 2023年年度报告 --cutoff 2024-01-31
python scripts/anchor_cli.py ask "中际旭创2024年年报增速区间" \
    --company 中际旭创 --ticker 300308.SZ \
    --period 2024年年度报告 --cutoff 2025-01-31
python scripts/anchor_cli.py audit <run_id>
\end{lstlisting}

{\sloppy
MCP服务提供五个工具。\texttt{anchor\_ask}对自然语言问题做完整裁决；\texttt{anchor\_case}运行具名案例；\texttt{anchor\_search}检索带日期的证据；\texttt{anchor\_audit}读取审计报告；\texttt{anchor\_contract}只解析断言契约，便于用户先检查阈值和结果空间。Claude Code通过项目根目录的\texttt{.mcp.json}连接服务，用户可以直接说“用anchor\_case运行zhongji\_2026，并解释每个断言的证据链”。\par
}

接口的安全边界如下。API密钥只从环境变量读取，不写入代码、日志或审计文件。下载的网页和PDF作为不可信文本处理，不执行其中的脚本。LLM响应按请求负载哈希缓存，缓存目录不进入公开审计文件。审计文件保留证据链接和原文引文，便于复核，但不包含API密钥或用户隐私。

接口的局限如下。网页检索结果会随时间和搜索引擎状态变化，不同运行日的证据集合可能不同。审计文件记录当次实际使用的证据，不能把某一次运行的后验当作固定输出。案例文件中的阈值由用户或案例作者给定，阈值不合理时，框架只能在给定契约内做一致的推断。接口是研究框架的演示，不是交易系统，也不构成投资建议。

